\documentclass[10pt,a4paper]{article}
\usepackage{amsmath,amssymb,amsthm,graphicx}
\title{SOS degree lower bound for exact planted clique}
\author{Shuo Pang}
\date{doi:10.4230/LIPIcs.CCC.2021.26; CC BY4.0}
\begin{document}
\maketitle
\noindent\textit{Editorial note.} Original author: Shuo Pang. Official CCC2021 article26 (2021-07-08), DOI10.4230/LIPIcs.CCC.2021.26.\par\medskip
\noindent\textit{Editorial note.} License: https://creativecommons.org/licenses/by/4.0/\par\medskip
\noindent\textit{Editorial note.} Faithful editable TeX collection of formal Theorem1.4 and the complete author source-local exact-clique lower-bound chain: exact model/parameters; required Section3 pseudoexpectation construction; Section4 homogeneity/concentration proofs and graphical-matrix definitions; Section5.1 Johnson calculation and Section5.3 full recursive factorization/matching/error proof; all Section6 Hadamard/Euler, exact factorization, factorial-Hankel structural positivity, pseudorandom/rest norm bounds and final PSD proof. Repeated statements are not duplicated. Standard cited prior norm/separator facts are exact prerequisites, not recursively re-proved. Optional non-exact mod-order by-product Section5.2/AppendixA and failed example3.9 are not targets or needed proof; the necessary Fact5.7 is retained. All author mathematical expressions/conditions remain uncorrected; no new proof or hypotheses. Complete published bibliography retained. Only official CC BY4.0 content used.\par\medskip
\noindent\textit{Allowed standard prerequisites (editorial).}\begin{itemize}
\item Graphical-matrix norm theorem. Exact author Theorem4.8[24,5], including its parameter bounds and stated moreover derivation, retained.
\item Min-separator lattice and standard graph separation. Exact author Fact5.7[11], Menger theorem and credited prior Lemma5.23[5] retained; the author own matching/factorization proofs remain complete.
\item Johnson schemes and SOS pseudoexpectation duality. Author model Definitions1.1/1.2, exact Johnson Fact5.3[9] and all own Johnson calculations retained.
\end{itemize}\medskip
\section*{Author exact-clique model and SOS conventions}
The SOS proof system, also known as the Positivstellensatz system [14], works with polynomials over $\mathbb R$. Given polynomial equalities (axioms) $f_1(\mathbf x)=0,\ldots,f_k(\mathbf x)=0$ on $\mathbf x=(x_1,\ldots,x_n)$, a proof (a refutation of the existence of a solution) is
\[-1=\sum_{i=1}^kf_iq_i+\sum_jr_j^2\quad\text{in }\mathbb R[x_1,\ldots,x_n],\]
where $q_1,\ldots,q_m$ and $r_1,\ldots$ are arbitrary polynomials on $x_1,\ldots,x_n$ over $\mathbb R$. Under certain conditions, in particular when all variables are boolean ($x_i^2=x_i$), such an refutation always exists if the axioms are contradictory. The degree-$d$ SOS proof system adds the limitation
\[\max_{i,j}\{\deg(f_i)+\deg(q_i),2\deg(r_j)\}\le d.\]
\paragraph{Definition 1.1.} Given an $n$-vertex simple graph $G$ and a number $\omega$, the Clique Problem for degree-$d$ SOS proof system has the following axioms:
\begin{align*}(\text{Boolean})\quad x_i^2&=x_i&&\forall i\in[n],\\(\text{Clique})\quad x_ix_j&=0&&\forall\{i,j\}\text{ non-edge},\\(\text{Size})\quad x_1+\cdots+x_n&=\omega.&&\tag{1.1}\end{align*}
To confirm no $\omega$-clique exists is to give an SOS refutation of the above. The SOS system has duality: to show degree lower bound it suffices to consider pseudo-expectation and the resulting moment matrix.\footnote{The author uses the name for simplicity. More cautiously, it should be called the pseudo-moment matrix.} With boolean variables this can be demonstrated on multi-linear polynomials. Let $X^{\le a}=\{x_S:S\subseteq[n],|S|\le a\}$ for any $a$.
\paragraph{Definition 1.2.} A degree-$d$ pseudo-expectation for the Clique Problem on $G$ is a map $\widetilde{\mathbb E}:X^d\to\mathbb R$ satisfying the following four constraints when extended by $\mathbb R$-linearity:
\begin{align*}(\text{Default})\quad\widetilde{\mathbb E}x_\emptyset&=1,\tag{1.2}\\(\text{Clique})\quad\widetilde{\mathbb E}x_S&=0&&\forall S:|S|\le d,\ G|_S\text{ non-clique},\tag{1.3}\\(\text{Size})\quad\widetilde{\mathbb E}((x_1+\cdots+x_n)x_S)&=\omega\widetilde{\mathbb E}x_S&&\forall S:|S|\le d-1.\tag{1.4}\end{align*}
In (1.4), $x_Ax_B:=x_{A\cup B}$. Define the moment matrix $M$ to be the $\binom{[n]}{\le d/2}\times\binom{[n]}{\le d/2}$ matrix with expression $M(A,B)=\widetilde{\mathbb E}x_{A\cup B}$ for all $|A|,|B|\le d/2$.\footnote{$d$ is always assumed to be even.} The last constraint is:\[M\text{ is positive semi-definite}.\tag{1.5}\]
If a degree-$d$ pseudo-expectation exists then there is no degree-$d$ SOS refutation. Exact Clique is Definition1.1. The Non-Exact Clique relaxation replaces the Size Constraints by one weaker inequality $\widetilde{\mathbb E}(x_1+\cdots+x_n)\ge\omega$. There is no planted clique in the problem's formulation; traditionally it is still called the planted clique problem due to the algorithmic motivation.

\section*{Formal author main result}
\paragraph{Theorem 1.4.} Let $\epsilon>0$ be any parameter, $\omega=n^{1/2-\epsilon}$. W.p. $>1-n^{-4\log n}$ over $G\sim G(n,\frac12)$, any SOS refutation of Exact Clique requires degree at least $\epsilon'\log n/\log\log n$, where $\epsilon'=\min\{\epsilon^2,1/40^2\}/2000$.

\section*{Author notation, parameter regime and Section 3}
$I,J,A,B,S$ denote vertex-sets and $T$ edge-sets. $E(S):=\binom S2$. $G$ denotes a simple graph on vertex-set $[n]$. ``$T\subseteq E([n])$'' is omitted in summation when there is no confusion. $y(n)=O(x(n))$ means there is some absolute constant $c$ such that $y(n)\le cx(n)$ for all $n$.
Throughout the paper,
\[\epsilon=\text{any positive parameter (wlog }\epsilon<1/40),\quad\omega=n^{1/2-4\epsilon},\quad\tau=\frac\epsilon{200}\frac{\log n}{\log\log n},\quad d=\frac\epsilon{100}\tau.\]
\subsection*{3.1 Non-exact case: a new perspective}
Given a graph $G$ we can think of a degree-$d$ pseudo-expectation as assigning a number $\widetilde{\mathbb E}x_S$ to each subset $S\subseteq[n]$ of size $\le d$, so that the resulting vector $\widetilde{\mathbb E}x$ looks indistinguishable to the expectation resulted from the case when a random-$\omega$ clique is planted, from the view of degree-$d$ SOS. To make up such an assignment we first go beyond to slightly larger subsets of size $\tau$. We define an ``approximate distribution'' on size $\le\tau$-cliques in $G$ then use it to generate pseudo-expectation on all size $\le d$-subsets.
\subsubsection*{3.1.1 $\zeta$-function and M\"obius inversion}
Given $n$-vertex graph $G$, let $p(G)\in\mathbb R^{2^{[n]}}$ be the max-clique-indicator vector, then $q(G):=\zeta p(G)$ is a vector supported exactly on all cliques in $G$, where $\zeta$ is the $2^{[n]}\times2^{[n]}$ matrix
\[\zeta(A,B)=1\text{ iff }A\subseteq B,\qquad\forall A,B\subseteq[n].\tag{3.1}\]
In particular, if $G$ itself is a single clique then $q(G)$ is the clique-indicator. $\zeta_{a,b}$ denotes the submatrix of $\zeta$ on rows $\binom{[n]}{\le a}$ and columns $\binom{[n]}{\le b}$; use similar notation on all related vectors. Consider the plant-situation where $G$ is indeed a single random clique. Suppose its distribution is represented by a plant-distribution vector $p_{\mathrm{plant}}\in\mathbb R^{2^{[n]}}$. Let the output-expectation $q_{\mathrm{out}}$ be indicator-vector of cliques in $G$ in expectation. Then
\[q_{\mathrm{out}}=\zeta p_{\mathrm{plant}}.\tag{3.2}\]
We call such a pair $(p_{\mathrm{plant}},q_{\mathrm{out}})$ a plant-setting.
\paragraph{Definition 3.1 (Two plant-settings).} The exact plant-setting $(p_0,q_0)$ is:
\[p_0(S)=\binom n\omega^{-1}\text{ if }|S|=\omega\text{ and }0\text{ otherwise},\tag{3.3}\]
and
\[q_0(S)=(\zeta p_0)(S)=\frac{\binom{n-|S|}{\omega-|S|}}{\binom n\omega}.\tag{3.4}\]
I.e. a random size-$\omega$ subset is chosen to be the planted clique. The independent plant-setting $(p_1,q_1)$ is:
\[p_1(S)=(\omega/n)^{|S|}(1-\omega/n)^{n-|S|}\quad\forall S\subseteq[n],\tag{3.5}\]
and
\[q_1(S)=(\zeta p_1)(S)=(\omega/n)^{|S|}.\tag{3.6}\]
I.e. any vertex is included in the planted clique w.p. $\omega/n$ independently.
Thus $\zeta$ reveals the basic linear relations between $(p_{\mathrm{plant}},q_{\mathrm{out}})$. It is upper-triangular (row- and column-indices ordered in a size-ascending way), invertible, and the inverse is the M\"obius inversion matrix:
\[\zeta^{-1}(A,B)=(-1)^{|B\setminus A|}\text{ if }A\subseteq B,\text{ and }0\text{ otherwise}.\]
Note $(\zeta_{a,a})^{-1}=(\zeta^{-1})_{a,a}$ for all $a\le n$. Moreover, if let the pseudo-expectation be defined as $\widetilde{\mathbb E}x=p\in\mathbb R^{2^{[n]}}$ for some vector $p$, then the ``full'' $2^{[n]}\times2^{[n]}$ moment matrix is
\[M_{\mathrm{SOS}}=\zeta\operatorname{diag}(p)\zeta^\top.\tag{3.7}\]
In particular, if $p$ is a nonnegative vector then $M_{\mathrm{SOS}}$ is immediately PSD.
\subsubsection*{3.1.2 The non-exact pseudo-expectation}
Given any $G$, first construct a degree-$\tau$ ``approximate plant-distribution'' $p_\tau(G)$, which simulates the plant-distribution (Definition3.1) in the sense that they give similar output-expectations. Require $p_\tau(G)$ to be supported on size $\le\tau$-cliques in $G$. Then take $\widetilde{\mathbb E}x=\zeta_{d,\tau}p_\tau(G)$ so the result inherits the linear structure posed by $\zeta$. From approximation it seems taking $\zeta_{\tau,\tau}^{-1}(q_1)_\tau$ would suffice, while to make it supported on cliques, same as in [13] add a clique-indicator factor:
\[p_\tau(G)(S)=2^{|\binom S2|}\mathrm{Cl}_S(G)\left(\zeta_{\tau,\tau}^{-1}(q_1)_\tau\right)(S)\quad\forall S\subseteq[n],\ |S|\le\tau.\tag{3.8}\]
Here $\mathrm{Cl}_S(\cdot)$ is the clique indicator function and $2^{|\binom S2|}$ is for re-normalization.
\paragraph{Definition 3.2.} $\forall S\subseteq[n]$, the normalized clique-indicator is function \[\widetilde{\mathrm{Cl}}_S(G):=2^{|\binom S2|}\mathrm{Cl}_S(G).\tag{3.9}\]
$\widetilde{\mathrm{Cl}}(G)$ denotes the column vector over a family of $S$'s, always clear from context.
\paragraph{Definition 3.3.} The non-exact pseudo-expectation is
\[\widetilde{\mathbb E}_{\mathrm{nonexact}}=\zeta_{d,\tau}p_\tau(G)=\zeta_{d,\tau}\left(\widetilde{\mathrm{Cl}}(G)\circ\zeta_{\tau,\tau}^{-1}\right)(q_1)_\tau\in\mathbb R^{\binom{[n]}{\le d}},\tag{3.10}\]
where $\circ$ is Hadamard product.\footnote{In general $(M_1\circ M_2)M_3\ne M_1\circ(M_2M_3)$, but they are equal if $M_1$ is a column vector.} In short, this refines the construction in [13] by one step: factor through size-$\tau$ subsets (in the only non-trivial way) so the size-$d$ output inherits linear relations posed by $\zeta$. The resulting moment matrix is
\[M_{\mathrm{nonexact}}(G)=\zeta_{d/2,\tau}\operatorname{diag}(p_\tau(G))(\zeta_{d/2,\tau})^\top,\tag{3.11}\]
similarly as (3.7).
\paragraph{Remark 3.4.} If $p_\tau(G)$ were nonnegative then PSDness would be immediate. Alas, this is not true by computation. PSDness could still possibly hold because $\zeta_{d/2,\tau}$ in (3.11) is degenerate.
\paragraph{Lemma 3.5 (Theorem1.5(1)).} For all $S\subseteq[n]$ with $|S|\le d$,
\[\widetilde{\mathbb E}_{\mathrm{nonexact}}x_S=\sum_{T:|V(T)\cup S|\le\tau}(\omega/n)^{|V(T)\cup S|}\chi_T.\tag{3.12}\]
\begin{proof}Note $\widetilde{\mathrm{Cl}}_S=\sum_{T\subseteq E(S)}\chi_T$ for all $S$. Now for $S,S'$ with appropriate size bound,
\[\left(\widetilde{\mathrm{Cl}}\circ\zeta_{\tau,\tau}^{-1}\right)(S,S')=\begin{cases}\sum_{T\in E(S)}\chi_T(-1)^{|S'\setminus S|},&S\subseteq S',\\0,&\text{o.w.}\end{cases}\]
Then
\begin{align*}&\left(\zeta_{d,\tau}(\widetilde{\mathrm{Cl}}\circ\zeta_{\tau,\tau}^{-1})\right)(S,S')\\&=\sum_{S'':S\subseteq S''\subseteq S'}\left(\sum_{T\subseteq E(S'')}\chi_T(-1)^{|S'\setminus S''|}\right)\\&=\sum_{T:V(T)\cup S\subseteq S'}\chi_T\left(\sum_{S'':V(T)\cup S\subseteq S''\subseteq S'}(-1)^{|S'\setminus S''|}\right)\\&=\sum_{T:V(T)\cup S\subseteq S'}\chi_T\delta_{S'=V(T)\cup S}\\&=\sum_{T:V(T)\cup S=S'}\chi_T.\end{align*}
Therefore,
\begin{align*}\widetilde{\mathbb E}_{\mathrm{nonexact}}x_S&=\left(\zeta_{d,\tau}(\widetilde{\mathrm{Cl}}\circ\zeta_{\tau,\tau}^{-1})(q_1)_\tau\right)(S)\\&=\sum_{S':|S'|\le\tau}\left(\sum_{T:V(T)\cup S=S'}\chi_T(\omega/n)^{|S'|}\right)\\&=\sum_{T:|V(T)\cup S|\le\tau}\chi_T(\omega/n)^{|V(T)\cup S|},\end{align*}
for all $S$ with $|S|\le d$.\end{proof}
\subsection*{3.2 The exact case}
In this subsection, we give a generic way to generate possible pseudo-expectations that satisfy Size Constraints (1.4). Define $\widetilde{\mathbb E}x_S$ in a top-down fashion: fix it for all $|S|=d$ first, then recursively set
\[\widetilde{\mathbb E}x_S\leftarrow\frac1{\omega-|S|}\sum_{i\notin S}\widetilde{\mathbb E}x_{S\cup\{i\}}\tag{3.13}\]
for smaller-sized $S$'s. Denote by $\widetilde{\mathbb E}_dx$ the assignments for $|S|=d$.
\paragraph{Definition 3.6.} The $d$-filtration matrix $\mathrm{Fil}_{d,=d}$, of dimension $\binom{[n]}{\le d}\times\binom{[n]}d$, is:
\[\mathrm{Fil}_{d,=d}(A,B)=\begin{cases}\binom{\omega-|A|}{d-|A|}^{-1},&A\subseteq B\ (|B|=d),\\0,&\text{otherwise}.\end{cases}\tag{3.14}\]
\paragraph{Definition 3.7.} Given vector $\widetilde{\mathbb E}_dx$ assigning a value to each $d$-subset $S\subseteq[n]$, the exact pseudo-expectation it generates is
\[\widetilde{\mathbb E}x:=\mathrm{Fil}_{d,=d}\widetilde{\mathbb E}_dx.\tag{3.15}\]
\paragraph{Lemma 3.8.} The pseudo-expectation in Definition3.7 satisfies the Size Constraints (1.4), regardless of the choice of $\widetilde{\mathbb E}_dx$.
\begin{proof}For any $S\in\binom{[n]}{<d}$, take $v_S\in\mathbb R^{\binom{[n]}{\le d}}$,
\[v_S(S')=\begin{cases}\omega-|S|,&S'=S,\\-1,&S'\supseteq S,\ |S'\setminus S|=1,\\0,&\text{otherwise}.
\end{cases}\]
It suffices to show $v_S^\top\mathrm{Fil}_{d,=d}=0$. But this is a direct check.\end{proof}
The generated $\widetilde{\mathbb E}$ should further satisfy Clique, PSDness and Default Constraints. Default is not a problem as long as $\widetilde{\mathbb E}x_\emptyset>0$, since we can rescale everything by $(\widetilde{\mathbb E}x_\emptyset)^{-1}$ without affecting other constraints.
\subsection*{3.3 The exact pseudo-expectation}
Take the pseudo-expectation for $|S|=d$ in the form
\[\widetilde{\mathbb E}x_S=\sum_{T:|V(T)\cup S|\le\tau}\chi_TF(|V(T)\cup S|),\]
for a function $F$ called a $d$-generating function. For general $|S|\le d$, (3.14) gives, putting $u:=d-|S|$,
\begin{align*}\widetilde{\mathbb E}x_S&=\frac1{\binom{\omega-d+u}u}\sum_{T:|V(T)\cup S|\le\tau}\chi_T\left[\sum_{c=0}^u\binom{|V(T)\cup S|-d+u}c\right.\\&\qquad\left.\cdot\binom{n-|V(T)\cup S|}{u-c}F(|V(T)\cup S|+u-c)\right].\tag{3.16}\end{align*}
\paragraph{Lemma 3.10.} Any exact pseudo-expectation generated by (3.16) satisfies Clique and Size Constraints (1.3),(1.4).
\begin{proof}For Clique Constraints, note (3.16) only depends on $|V(T)\cup S|$, so by grouping terms $\widetilde{\mathbb E}x_S=\sum_{T:|V(T)\cup I\cup J|\le\tau}M(I,J;T)\chi_T$ factors through $\widetilde{\mathrm{Cl}}_{I\cup J}=\sum_{T\subseteq E(I\cup J)}\chi_T$. I.e., $M(I,J)(G)=0$ if $\widetilde{\mathrm{Cl}}_{I\cup J}(G)=0$. It satisfies Size Constraints by Lemma3.8.\end{proof}
\paragraph{Definition 3.11 (Exact $d$-generating function).}
\[F(x):=\frac{(x+8\tau^2)!}{(8\tau^2)!}(\omega/n)^x.\]
\paragraph{Remark 3.12.} The design of $F$, especially its first factor, is technical and the ultimate goal is to make the resulting $M$ positive. The numerator $(x+8\tau^2)!$ will be used in Proposition6.28, where $8\tau^2$ can be replaced by larger polynomials in $\tau$. The $(8\tau^2)!$ denominator is just for convenience (Remark3.14).
\paragraph{Definition 3.13.} The exact moment matrix $\widetilde M$ is
\[\widetilde M(A,B)=\sum_{T:|V(T)\cup A\cup B|\le\tau}\widetilde M(A,B;T)\chi_T\quad\forall |A|,|B|\le d/2,\]
where, putting $u=d-|A\cup B|$,
\begin{align*}\widetilde M(A,B;T)&=\frac1{\binom{\omega-d+u}u}\left[\sum_{c=0}^u\binom{|V(T)\cup A\cup B|-(d-u)}c\right.\\&\qquad\cdot\binom{n-|V(T)\cup A\cup B|}{u-c}\\&\qquad\left.\cdot\frac{(|V(T)\cup A\cup B|+u-c+8\tau^2)!}{(8\tau^2)!}(\omega/n)^{|(V(T)\cup A\cup B)|+u-c}\right].\tag{3.17}\end{align*}

\section*{Author Section 4: preparations}
\paragraph{Fourier notation.} $\chi_T$ is the Fourier character for the edge-set $T$ in the random graph. The author's Fourier orthonormal basis in Equation(7.1) is $\chi_T=\prod_{e\in T}(x_e-(2p-1))/(2\sqrt{p(1-p)})$ for $G(n,p)$; at $p=1/2$ this is the product of the independent signed edge indicators. This is the standard Fourier convention used throughout the paper.
\paragraph{Remark 3.14 (continued).} In (3.17), the ``most significant'' factor is $(\omega/n)^{|V(T)\cup A\cup B|}\omega^{-c}$, if notice
\[\frac{\binom{n-|V(T)\cup A\cup B|}{u-c}}{\binom{\omega-d+u}u}\omega^un^{-(u-c)}\]
has 0th-order in $\omega,n$. Other factors like $(|V(T)\cup A\cup B|+u-c+8\tau^2)!/(8\tau^2)!$ are qualitatively smaller than $\omega$, within our parameter regime.
\subsection*{4.1 Homogenization for Exact Clique}
With Size Constraints (1.4) satisfied, any moment matrix can be reduced to its $\binom{[n]}{d/2}$-principal minor, which is slightly more convenient. The homogeneity trick is standard in SOS literature. Given a degree-$d$ moment matrix $M^d_{\mathrm{SOS}}(G)$ satisfying (1.4), let $M(G)$ be its principal minor on $\binom{[n]}{d/2}\times\binom{[n]}{d/2}$.
\paragraph{Lemma 4.1.} $M^d_{\mathrm{SOS}}(G)$ is PSD iff $M(G)$ is PSD.
\begin{proof}The $\Rightarrow$ part is trivial. Now suppose $M^d_{\mathrm{SOS}}$ is not PSD, then
\[\exists a\in\mathbb R^{\binom{[n]}{\le d/2}}\quad a^\top M^d_{\mathrm{SOS}}a=-1.\tag{4.1}\]
With boolean constraints (define $\widetilde{\mathbb E}(x_i^2p):=\widetilde{\mathbb E}(x_ip)$ for all $i$ and polynomial $p$ of degree $\le d-2$), this is equivalent to $\widetilde{\mathbb E}(g^2)=-1$ (4.2), where $g=a^\top x=\sum_{|S|\le d/2}a_Sx_S$ is multi-linear. Now substitute every $x_S$ ($|S|<d$) in $g$ by the corresponding linear combination of $\{x_{S'}:|S'|=d\}$ from (3.13). This does not affect (4.2) since $\widetilde{\mathbb E}$ satisfies the equality constraints. We get $\widetilde{\mathbb E}(g_1^2)=-1$ (4.3) for some multi-linear, degree-$d/2$ homogeneous $g_1$. Translate (4.3) back (assume $g_1=b^\top x$, $x=(x_S)_{|S|=d/2}$) to $b^\top Mb=-1$, so $M$ is not PSD.\end{proof}
\subsection*{4.2 Concentration bound on polynomials}
\paragraph{Lemma 4.2.} Suppose $a<\log n$, and $p=\sum_{T:|V(T)|=a}c(T)\chi_T$, $c_T\in\mathbb R$, is a polynomial, and $C>0$ satisfies $|c(T)|\le C$ for all $T$. Then w.p. $1-n^{-10\log n}$ over $G$,
\[|p(G)|<C\,n^{a/2}2^{a^2}n^{4\log\log n}.\tag{4.4}\]
\begin{proof}Power-estimation. For all $k\in\mathbb N$ (think of $a<k=o(n/a)$),
\[p^{2k}=\sum_{T_1,\ldots,T_{2k}:|V(T_i)|=a}c(T_1)\cdots c(T_{2k})\chi_{T_1}\cdots\chi_{T_{2k}}.\tag{4.5}\]
Take expectation. Each $\mathbb E[\chi_{T_1}\cdots\chi_{T_{2k}}(G)]\ne0$ (equals 1) iff every edge appears even times in $T_1,\ldots,T_{2k}$, which implies $|V(T_1\cup\cdots\cup T_{2k})|\le\frac12\cdot2ka=ka$. There are at most $\binom n{ka}<n^{ka}$ choices of such $V(T_1\cup\cdots\cup T_{2k})$. For each choice, at most $\binom{ka}a2^{\binom a2}<(ka)^a2^{a^2/2}$ ways to choose each $T_i$. Therefore,
\[\mathbb E[p^{2k}]\le C^{2k}n^{ka}((ka)^a2^{a^2/2})^{2k}=:N^{2k},\quad N=Cn^{a/2}(ka)^a2^{a^2/2}.\]
By Markov, $\Pr[p^{2k}>(2N)^{2k}]<2^{-2k}$. Take $k:=10\log^2n$; w.p. $>1-n^{-10\log n}$, $|p(G)|<2N<Cn^{a/2}2^{a^2}n^{4\log\log n}$ for all large enough $n$.\end{proof}
\subsection*{4.3 Norm concentration of pseudo-random matrices}
Pseudo-random matrices refer to graphical matrices [24]: intuitively, such a matrix collects Fourier characters of all embeddings of a fixed ``shape''. Definitions4.3/4.5 are implicit in [24,25,16] and termed explicitly in [5].
\paragraph{Definition 4.3.} A ribbon $R$ is an ordered triple $(A,B;T)$ with $A,B$ vertex-sets and $T$ an edge-set. $A,B$ are its left/right vertex-sets. Its size is $|V(R)|=|V(T)\cup A\cup B|$. By definition a ribbon as a graph has no isolated vertex outside $A\cup B$.
\paragraph{Definition 4.4.} $R=(A,B;T)$ is left-generated if every vertex in $V(R)$ is either in $B$ or can be reached by paths from $A$ without touching $B$.\footnote{The author always includes degenerate paths (one-point paths).} Right-generated is symmetric.
\paragraph{Definition 4.5.} A ribbon $(A,B;T)$ is a shape if $A\cup B$ is totally ordered. Denote a shape by $U=(A,B;T)$, with $V(U)=A\cup B\cup V(T)$, size $|V(U)|$. When fixing underlying vertex-set $[n]$, a ribbon within $[n]$ can always be regarded as shapes with the induced vertex ordering. The author may interchangeably use $R$ to denote shapes.
\paragraph{Definition 4.6.} A real-valued function $f$ on ribbons within $[n]$ is symmetric with respect to shapes if $f(R)=f(R')$ whenever $R,R'$ have the same shape.
\paragraph{Definition 4.7 ([24]).} Fix $n$ and shape $U=(A,B;T)$. A map $\phi:V(U)\to[n]$ is proper if injective and respects the order on $A\cup B$. Its graphical matrix is the $2^{[n]}\times2^{[n]}$ matrix
\[\forall I,J\subseteq[n],\quad M_U(I,J)=\sum_{T':\exists\text{ proper }\phi,\ \phi(A)=I,\phi(B)=J,\phi(T)=T'}\chi_{T'}\]
(0 if no such $\phi$ exists). Here $\phi$ on $T$ is the natural induced edge map.
\paragraph{Theorem 4.8 (Norm bounds on $M_U$ [24,5]).} For any shape $U=(A,B;T)$ of size $t<\log n$, w.p. $>1-n^{-10\log n}$ over $G$,
\[\|M_U(G)\|\le n^{(t-p)/2}2^{O(t)}(\log n)^{O(t+p-2r)},\tag{4.6}\]
where $r=|A\cap B|$, $p$ is the maximum number of vertex-disjoint paths between $(A,B)$ in $U$. This is tight up to $\mathrm{polylog}(n)$ factors for all $M_U$ with these parameters ([24],Theorem38). Moreover, if $s=(|A|+|B|)/2$, then
\[\|M_U(G)\|\le n^{(t-p)/2}2^{O(t)}(\log n)^{O(t-s)}.\tag{4.7}\]
The author cites the standard trace-power norm theorem without reproducing its earlier proof. The ``moreover'' part follows since $t\ge|A\cup B|=2s-r$, $p\le s$, so $t+p-2r\le t+s-2(2s-t)=3(t-s)$.
\paragraph{Remark 4.9.} Shapes $U=(A,A;T)$ give diagonal $M_U$, so estimating its norm is estimating the absolute values of its polynomial diagonals; Theorem4.8 generalizes Lemma4.2.
\subsection*{4.4 General notions on graphs}
\paragraph{Definition 4.10 (Vertex-separator).} Given graph $H$ and $A,B\subseteq V(H)$, call $S\subseteq V(H)$ an $(A,B)$-vertex-separator if every path from $A$ to $B$ in $H$ must pass through $S$.\footnote{Degenerate paths are included, as before.} Let $s_{A,B}(H):=\min\{|S|:S\text{ is an }(A,B)\text{-vertex-separator}\}$. A separator achieving this is a min-separator; $\mathrm{mSep}_{A,B}(H)$ is the set of all min-separators. For ribbon $R=(A,B;T)$ use the graph on $V(T)\cup A\cup B$ with edges $T$ and write $s_{A,B}(T)$, $\mathrm{mSep}_{A,B}(T)$ or $\mathrm{mSep}(R)$.
\paragraph{Menger's theorem.} For finite graph $H$, $s_{A,B}(H)$ equals the maximum number of vertex-disjoint paths from $A$ to $B$ in $H$.
\paragraph{Definition 4.11.} The reduced size of ribbon $R=(A,B;T)$ is\[e_{A,B}(T):=|V(T)\cup A\cup B|-s_{A,B}(T).\tag{4.8}\]
It is double the exponent of $n$ in Theorem4.8's bound, hence controls the norm of the graphical matrix.

\section*{Required author Section 5.1: diagonalization of the expectation}
Throughout Section5, $M'$ denotes the $\binom{[n]}{d/2}\times\binom{[n]}{d/2}$ minor of the non-exact moment matrix:
\[M'(I,J)=\sum_{T:|V(T)\cup I\cup J|\le\tau}(\omega/n)^{|V(T)\cup I\cup J|}\chi_T,\quad\forall I,J\in\binom{[n]}{d/2}.\tag{5.1}\]
\footnote{Strictly speaking PSDness of this minor is not sufficient since there is no homogeneity reduction in the non-exact case; it suffices to demonstrate the idea.} The goal is to diagonalize $M'$ approximately so that the difference matrix is negligible w.h.p. on plugging $G$.
\paragraph{Proposition 5.1.} $\mathbb E[M']=CC^\top$, where $C$ is the $\binom{[n]}{d/2}\times\binom{[n]}{\le d/2}$ matrix
\[C=(\zeta^\top)_{d/2,\le d/2}\operatorname{diag}(\sqrt{t(|A|)})_{A\in\binom{[n]}{\le d/2}},\tag{5.2}\]
and $t(r)=(1-O(d\omega/n))(\omega/n)^{d-r}$ for $r=0,\ldots,d/2$.
\paragraph{Definition 5.2 (See e.g. [9]).} Fix $n,k$. A Johnson scheme $J$ is a $\binom{[n]}k\times\binom{[n]}k$ matrix satisfying $J(I,J)=J(I',J')$ whenever $|I\cap J|=|I'\cap J'|$. For fixed $n,k$, Johnson schemes are symmetric matrices and form a commutative $\mathbb R$-algebra, hence are simultaneously diagonalizable. Fix $k=d/2$. One basis is $D_0,\ldots,D_{d/2}$ where
\[D_r(I,J)=\begin{cases}1,&|I\cap J|=r,\\0,&\text{o.w.}\end{cases}\tag{5.3}\]
Another basis is $J_0,\ldots,J_{d/2}$,
\[J_r(I,J)=\binom{|I\cap J|}r.\tag{5.4}\]
These are PSD since
\[J_r=\sum_{A\subseteq[n],|A|=r}u_Au_A^\top,\quad u_A\in\mathbb R^{\binom{[n]}k},\ u_A(B)=\mathbf1_{A\subseteq B}.\tag{5.5}\]
Clearly $J_{d/2}=\mathrm{Id}$.
\paragraph{Fact 5.3 (See e.g. (4.29) in [9]).} The Johnson schemes for $(n,d/2)$ have shared eigenspace-decomposition $\mathbb R^{\binom{[n]}{d/2}}=V_0\oplus\cdots\oplus V_{d/2}$, and $J_r=\bigoplus_{i=0}^{d/2}\lambda_r(i)\Pi_i$, where $\Pi_i$ is orthogonal projection to $V_i$ w.r.t. Euclidean inner product, and
\[\lambda_r(i)=\binom{d/2-i}{r-i}\binom{n-d/2-i}{d/2-r},\qquad0\le i\le d/2.\]
\paragraph{Lemma 5.4.} $\mathbb E[M']=\sum_{r=0}^{d/2}t(r)J_r$, with $t(r)=(1-O(d\omega/n))(\omega/n)^{d-r}$.
\begin{proof}By definition $\mathbb E[M']=\sum_{r=0}^{d/2}(\omega/n)^{d-r}D_r$. Each $D_r$ decomposes as
\[D_r=\sum_{r'=r}^{d/2}(-1)^{r'-r}\binom{r'}rJ_{r'}.\tag{5.6}\]
Indeed the RHS at $(I,J)$ equals
\begin{align*}&\sum_{r'=r}^{d/2}(-1)^{r'-r}\binom{r'}r\binom{|I\cap J|}{r'}\\&=\sum_{r'=r}^{|I\cap J|}(-1)^{r'-r}\binom{|I\cap J|}r\binom{|I\cap J|-r}{r'-r}\\&=\binom{|I\cap J|}r\mathbf1_{|I\cap J|=r}=\mathbf1_{|I\cap J|=r}.\end{align*}
So
\begin{align*}\mathbb E[M']&=\sum_{r=0}^{d/2}(\omega/n)^{d-r}\left(\sum_{r'=r}^{d/2}(-1)^{r'-r}\binom{r'}rJ_{r'}\right)\\&=\sum_{r'=0}^{d/2}J_{r'}\left(\sum_{r=0}^{r'}(\omega/n)^{d-r}(-1)^{r'-r}\binom{r'}r\right)\\&=\sum_{r'=0}^{d/2}J_{r'}(\omega/n)^{d-r'}(1-\omega/n)^{r'}.\tag{5.7}\end{align*}
This proves the lemma.\end{proof}
By Lemma5.4 and (5.5), letting $t(r')=(\omega/n)^{d-r'}(1-\omega/n)^{r'}$, we have
\begin{align*}\mathbb E[M']&=\sum_{A:|A|\le d/2}t(|A|)u_Au_A^\top\\&=(\zeta^\top)_{d/2,\le d/2}\operatorname{diag}(t(|A|))\zeta_{\le d/2,d/2}=CC^\top,\end{align*}
since $(\zeta^\top)_{d/2,\le d/2}$ has columns $\{u_A:|A|\le d/2\}$. This proves Proposition5.1.
\paragraph{Fact 5.7 ([11], required separator notation from Section5.2).} For ribbon $(I,J;T)$, all min-separators $\mathrm{mSep}_{I,J}(T)$ have a natural poset structure: $A_1\le A_2$ iff $A_1$ separates $(I,A_2;T)$, equivalently iff $A_2$ separates $(J,A_1;T)$. This poset is a lattice: each two min-separators have join and meet. In particular it has a unique minimum and maximum, denoted $S_l(I,J;T)$ and $S_r(I,J;T)$, the leftmost and rightmost min-separator. The rest of Section5.2 concerns the optional non-exact by-product; the author expressly permits skipping that section for the exact-clique proof.

\section*{Required author Section 5.3: recursive factorization}
This formalizes and extends recursive factorization, giving notions extendable to matrix products (Definitions5.13/5.15), a simplification (Lemma5.25) and refinement (Proposition5.24) for later use.
\paragraph{Definition 5.9.} Let $L$ be the $\binom{[n]}{d/2}\times\binom{[n]}{\le d/2}$ matrix
\[L(I,A):=\sum_{\substack{T':|V(T')\cup I\cup A|\le\tau,\ A=S_l(I,A;T')\\T'\cap E(A)=\emptyset,\ (I,A;T')\text{ left-generated}}}(\omega/n)^{|V(T')\cup I\cup A|-|A|}\chi_{T'}.\tag{5.19}\]
Let $Q_0$ be the $\binom{[n]}{\le d/2}\times\binom{[n]}{\le d/2}$ matrix
\[Q_0(A,B):=\sum_{\substack{T_m:|T_m\cup A\cup B|\le\tau\\A,B\in\mathrm{mSep}_{A,B}(T_m)}}(\omega/n)^{|V(T_m)\cup A\cup B|}\chi_{T_m}.\tag{5.20}\]
Finally, $D:=\operatorname{diag}((\omega/n)^{|A|/2})_{A\in\binom{[n]}{\le d/2}}$ (5.21). The author calls $L(DQ_0)L^\top$ the first-approximate diagonalization of $M'$. Despite the name, $M'-L(DQ_0D)L^\top$ (5.22) is far from negligible. Recursive factorization will in the end give
\[M'=L[D(Q_0-Q_1+Q_2\cdots\pm Q_{d/2})D]L^\top+E\tag{5.23}\]
for negligible error-matrix $E$.
\paragraph{Remark 5.10.} Use of $D$ is superficial; it keeps middle matrices $Q_i$ better-positioned. $LD$ corresponds to the ``$L$'' matrix of [5].
\paragraph{Definition 5.11 ([5],Definition6.5).} The canonical decomposition of ribbon $R=(I,J;T)$ is the uniquely determined ribbon-triple
\[(R_l,R_m,R_r)=((I,A;T_l),(A,B;T_m),(B,J;T_r)),\]
where $A=S_l(I,J;T)$, $B=S_r(I,J;T)$. $V(R_l)$ is $A$ unioned with vertices reachable from $I$ in $T$ without touching $A$, and $T_l=T|_{V(R_l)}\setminus E(A)$. $V(R_r)$ is symmetric, reachable from $J$ without touching $B$, with $T_r=T|_{V(R_r)}\setminus E(B)$. Finally $T_m=T\setminus(T'\sqcup T'')$. These are the left, middle and right ribbon of $R$.
\paragraph{Remark 5.12.} The canonical decomposition satisfies
\begin{gather*}A=S_l(I,A;T_l),\quad B=S_r(B,J;T_r),\\T_l\cap E(A)=\emptyset=T_r\cap E[A],\\R_l\text{ left-generated},\quad R_r\text{ right-generated},\\A,B\in\mathrm{mSep}_{A,B}(T_m)\quad(\text{so }|A|=|B|).
\end{gather*}
These are the inner conditions. The outer intersection property is
\[V(R_l)\cap V(R_m)\subseteq A,\quad V(R_m)\cap V(R_r)\subseteq B,\quad V(R_l)\cap V(R_r)\subseteq A\cap B,\]
implying $e(R_l)+|V(R_m)|+e(R_r)=|V(R)|$ (5.24).
\paragraph{Definition 5.13 (Inner and outer canonicality).} For a triple of ribbons $(R_l,R_m,R_r)=((I,A;T_l),(A,B;T_m),(B,J;T_r))$ ($T_l,T_m,T_r$ arbitrary subsets of an edge-set), their ribbon-sum is $(I,J;T)$ where $T=T_l\oplus T_m\oplus T_r$. It is inner-canonical if it satisfies the four inner conditions above (5.25). It is outer-canonical if
\[V(R_l)\cap V(R_m)\subseteq A,\quad V(R_m)\cap V(R_r)\subseteq B,\quad V(R_l)\cap V(R_r)\subseteq A\cap B.\tag{5.26}\]
It is canonical if both inner- and outer-canonical.
\paragraph{Proposition 5.14.} Canonical triples are 1-1 correspondent to their ribbon-sum, via canonical decomposition.
\begin{proof}This follows by an immediate check from the definition.\end{proof}
Recall $\mathbb R[\{\chi_T\}]$ is the ring adding fresh variables $\chi_T$ into $\mathbb R$ for each $T\subseteq\binom{[n]}2$, with relations $\chi_{T'}\chi_{T''}=\chi_T$ whenever $T'\oplus T''=T$.
\paragraph{Definition 5.15 (Approximate form).} Suppose matrices $X,Y$ have rows/columns indexed by subsets of $[n]$, entries in $\mathbb R[\{\chi_T\}]$, and every character viewed as a ribbon on (row,column) has size $\le\tau$. Suppose $XYX^\top$ is defined. Every nonzero triple product, without collecting like terms, is
\[X(I,A;T_l)Y(A,B;T_m)X(J,B;T_r)\chi_{T_l}\chi_{T_m}\chi_{T_r}\tag{5.28}\]
in $XYX^\top$ (5.27), with the coefficient nonzero in $\mathbb R$. Identify it with the natural ribbon-triple; its resulting term is $X(I,A;T_l)Y(A,B;T_m)X(J,B;T_r)\chi_{T_l\oplus T_m\oplus T_r}$. Call it outer-canonical if that triple is outer-canonical; it exceeds degree if $|V(T)\cup I\cup J|>\tau$. The approximation form is
\begin{align*}XYX^\top&=(XYX^\top)_{\mathrm{can}}+(XYX^\top)_{\mathrm{non\text{-}can}}+E_{\mathrm{deg}},\\(XYX^\top)_{\mathrm{can}}&=XYX^\top-(XYX^\top)_{\mathrm{non\text{-}can}}-E_{\mathrm{deg}}.\tag{5.29}\end{align*}
Here the canonical part collects outer-canonical products not exceeding degree, the non-canonical part collects non-outer-canonical products, and $E_{\mathrm{deg}}$ all remaining terms.
\paragraph{Remark 5.16.} Proposition5.14 explains Definition5.9 a posteriori: $M'=[L(DQ_0D)L^\top]_{\mathrm{can}}$.
\paragraph{Definition 5.17 (First-approximate factorization of $M'$).}
\[M'=L(DQ_0D)L^\top-[L(DQ_0D)L^\top]_{\mathrm{non\text{-}can}}-E_{1;\mathrm{deg}}.\tag{5.30}\]
$E_{1;\mathrm{deg}}$ is Definition5.15 applied to this product. $L(DQ_0D)L^\top$ is the first-approximate factorization. The high-degree error is negligible in norm (proved analogously in the exact case); matrices considered are supported on clique rows/columns of $G$. The main error is the non-canonical part. The key is that it too approximately factors through $L,L^\top$: there is $Q_1$ with
\[[L(DQ_0D)L^\top]_{\mathrm{non\text{-}can}}=[L(DQ_1D)L^\top]_{\mathrm{can}}+E'_{1;\mathrm{negl}}.\tag{5.31}\]
Repeat for $[L(DQ_1D)L^\top]_{\mathrm{non\text{-}can}}$ and so on.
\paragraph{Definition 5.18 ([5],Definition6.9).} A generalized ribbon is a usual ribbon with a new isolated vertex-set. Write $R^*=(A,B;T^*)$ with $T^*=T\sqcup I$, $T$ an edge-set, $I$ disjoint from $V(T)\cup A\cup B$, denoted $I(R^*)$. $V(R^*)=V(T)\cup A\cup B\cup I$. Usual ribbons have $I=\emptyset$. $(A,B;T)$ is the unique largest ribbon in $R$.
\paragraph{Remark 5.19.} $I(R^*)$ differs from the isolated set of the underlying graph since it excludes $A\cup B$.
\paragraph{Definition 5.20.} A side-inner-canonical triple $((I,A;T_l),(A,B;T_m),(B,J;T_r))$ has left/right ribbons satisfying their inner-canonical conditions (first three of (5.25)); the middle is just a ribbon.
\paragraph{Definition 5.21 (Separating factorization; [5],Definition6.10).} Given a side-inner-canonical triple $(R_l,R_m,R_r)$ as above, put $T=T_l\oplus T_m\oplus T_r$. Let $Z$ be the multiset of unexpected intersections, the vertices from $(R_l\cap R_m)-A$, $(R_m\cap R_r)-B$, $(R_l\cap R_r)-(A\cap B)$. Its intersection size is $z(R_l,R_m,R_r)=|Z|$. Then
\[|V(R_l)\cup V(R_m)\cup V(R_r)|=|V(R_l)|+|V(R_m)|+|V(R_r)|-|A|-|B|-z.\tag{5.32}\]
Define $S'_l$ to be the leftmost min-separator of $(I,A\cup(Z\cap V(R_l));T_l)$, and $S'_r$ the rightmost min-separator of $(B\cup(Z\cap V(R_r)),J;T_r)$. They lie in $V(T)\cup I\cup J$. Define $R'_l=(I,S'_l;T'_l)$: $V(R'_l)$ is $S'_l$ unioned with vertices in $R_l$ reachable from $I$ by paths in $T_l$ without touching $S'_l$, and $T'_l$ is $T_l\setminus E(S'_l)$ restricted to $V(R'_l)$. Define $R'_r$ symmetrically. In particular $T'_l\cap T'_r=\emptyset$. The generalized middle ribbon is $R^*_m=(S'_l,S'_r;T^*_m)$ with
\[T^*_m=T\setminus(T'_l\sqcup T'_r)\sqcup I(R^*_m),\]
where $I(R^*_m)$ collects all remaining isolated vertices:
\[I(R^*_m)=V(R_l)\cup V(R_m)\cup V(R_r)-V(T)\cup I\cup J.\tag{5.33}\]
This separating factorization is denoted $(R_l,R_m,R_r)\to(R'_l,R^*_m,R'_r)$ (5.34).
\paragraph{Remark 5.22.} (1) The resulting triple is side-inner-canonical and outer-canonical. Its pairwise intersections lie within $S'_l,S'_r,S'_l\cap S'_r$. Replacing $R^*_m$ by its largest ribbon gives a canonical triple disjoint from $I(R^*_m)$. (2) $R'_l\subseteq R_l$, and $S'_l$ separates $(V(R'_l),V(R_l)-V(R'_l))$ in $R_l$. The part of $R_l$ to the right of $S'_l$ is disjoint from $R'_l$ and is in $R^*_m$; symmetrically for $R_r$. (3) Since $S'_l$ separates $(I,A)$ in $R_l$ and $A$ is its unique min-separator, there are $|A|$ vertex-disjoint paths from $A$ to $S'_l$ in $R_l$; symmetrically for $R_r$.

\section*{Required author factorization lemmas, Sections 5.3.2--5.3.3}
\paragraph{Lemma 5.23 ([5],Lemmas6.14/7.14).} Suppose $(R_l,R_m,R_r)\to(R'_l,R^*_m,R'_r)$. In Definition5.21's notation: (1) $|S'_l|+|S'_r|\ge|A|+|B|+1$. (2) Put $s=(|A|+|B|)/2$, let $p'$ be the maximum number of vertex-disjoint paths from $S'_l$ to $S'_r$ in $R^*_m$, and $p$ from $A$ to $B$ in $R_m$. Then
\[2(s'-s)+(p-p')+|I(R^*_m)|\le z(R_l,R_m,R_r).\]
\begin{proof}(1) There must be an unexpected pairwise intersection. In any of the three cases breaking (5.26), there is $v\in Z$ in $V(R_l)-A$ or $V(R_r)-B$. WLOG the first. Then $S'_l\ne A$ since by left-generatedness $v$ is reachable from $I$ without passing $A$. Similarly, if $|S'_l|=|A|$ it is $A$, since $A$ is the unique min-separator separating $(I,A)$, impossible. Thus $S'_l>A$. (2) The author refers to the proof in the cited original paper.\end{proof}
\footnote{In this setting $R_m$ is a ribbon without isolated vertices. The credited prior lemma is a standard prerequisite, not a separately selected contribution.}
The separating factorization cancels a term in $[L(DQ_0D)L^\top]_{\mathrm{non\text{-}can}}$ by a canonical term using $R'_l$, $R'_r$ and the largest ribbon of $R^*_m$ in a new middle matrix $DQ_1D$. Other triples can have the same separating factorization, and entries of $L$ are sums of different left ribbons. Simultaneous cancellation requires the following refined solvability condition, also needed in Lemma6.20.
\paragraph{Proposition 5.24 (Solvability condition, cf.[5],Claim6.12).} Fix $(I,J,S'_l,S'_r)$ and generalized ribbon $R^*_m$ on $(S'_l,S'_r)$. Let $(R'_l,R'_r)$ be inner-canonical left/right ribbons on $(I,S'_l)$, $(S'_r,J)$. Let $(R''_l,R''_r)$ be another such pair with $e(R'_l)=e(R''_l)$, $e(R'_r)=e(R''_r)$ (equivalently same sizes). For every fixed $(i,j,z)$ there is a 1-1 matching between
\[\{(R_l,R_m,R_r): (R_l,R_m,R_r)\to(R'_l,R^*_m,R'_r),\ (e(R_l),e(R_r),z(R_l,R_m,R_r))=(i,j,z)\}\tag{5.35}\]
and the analogous set with $(R''_l,R^*_m,R''_r)$ (5.36). The matching fixes each middle $R_m$.
\begin{proof}Give a reversible map from (5.35) to (5.36). Take a triple in (5.35). By Remark5.22(2), the part of $R_l$ right of $S'_l$ is in $R^*_m$, hence disjoint from $R'_l$ and $R''_l$; symmetrically for the right side. Map
\[(R_l,R_m,R_r)\mapsto(\phi(R_l),R_m,\phi(R_r)),\]
where $\phi(R_l)$ replaces $R'_l$ by $R''_l$ in $R_l$, and symmetrically on the right. $R^*_m$, hence $R_m$, is unchanged. Since $R'_l,R''_l$ have the same size and the above parts are disjoint, this keeps $|V(R_l)|$. $R_l,\phi(R_l)$ have the same right distinguished set, the unique min-separator of both, so $e(R_l)=e(\phi(R_l))$; symmetrically on the right. Thus $(i,j)$ is unchanged. $z$ is unchanged since the changed part is disjoint from $Z(R_l,R_m,R_r)$. Reverse by switching the roles of the two ribbon pairs.\end{proof}
\paragraph{Lemma 5.25 (One round of factorization).} Let $L$ be (5.19) and $Q$ any $\binom{[n]}{\le d/2}\times\binom{[n]}{\le d/2}$ matrix with entries
\[Q(A,B)=\sum_{T_m:|V(T_m)\cup A\cup B|\le\tau}(\omega/n)^{|V(R_m)|}q(R_m)\chi_{T_m},\tag{5.37}\]
where $R_m=(A,B;T_m)$ and $q$ is symmetric w.r.t. shapes. Define $Q'$ and $E'_{\mathrm{negl}}$ so that
\[(LQL^\top)_{\mathrm{non\text{-}can}}=(LQ'L^\top)_{\mathrm{can}}+E'_{\mathrm{negl}}.\tag{5.38}\]
First set
\[Q'(A,B)=\sum_{T_m:|V(T_m)\cup A\cup B|\le\tau}(\omega/n)^{|V(R_m)|}q'(R_m)\chi_{T_m}.\tag{5.39}\]
Fix $R_m=(A,B;T_m)$ with $t=|V(R_m)|\le\tau$, $s=(|A|+|B|)/2$. For each generalized ribbon $R^*_m$ whose largest ribbon is $R_m$ and $|V(R^*_m)|\le\tau$, fix $(R'_l,R'_r)$ such that $(R'_l,R^*_m,R'_r)$ is the separating factorization of some triple with $|V(R'_l)|,|V(R'_r)|\le\tau$; exclude $R^*_m$ if none. Let
\begin{align*}q'(R_m)&=\sum_{\substack{R^*_m:\text{generalized on }(A,B)\\|V(R^*_m)|\le\tau,\ \text{largest ribbon }R_m}}(\omega/n)^{|I(R^*_m)|}q''(R^*_m),\\q''(R^*_m)&=\sum_{1\le z\le d/2}\sum_{\substack{\mathcal P=(R_l,R,R_r):\text{ side-inner-canonical}\\\mathcal P\to(R'_l,R^*_m,R'_r),\ z(\mathcal P)=z}}(\omega/n)^zq(R).\tag{5.40}\end{align*}
By Proposition5.24 this is independent of the fixed pair and remains symmetric w.r.t. shapes. Define $E'_{\mathrm{negl}}$ by (5.38). Conclusions: (1) w.p. $>1-n^{-9\log n}$,
\[\|E'_{\mathrm{negl}}\|\le\max\{q(A,B;T)\}\,n^{-\epsilon\tau}.\]
(2) If $C$ satisfies
\[\forall R_m\quad|q(R_m)|\le C(\omega/n^{1-\epsilon})^{s-p},\tag{5.41}\]
where $p$ is the maximum number of vertex-disjoint paths between $A,B$ in $R_m$ (equal to $s_{A,B}(T_m)$ by Menger), then
\[\forall R_m\quad|q'(R_m)|\le C(\omega/n^{1-\epsilon})^{s-p+1/3}.\]
\begin{proof}Compare the canonical and non-canonical parts as (0), then prove (1),(2).
\paragraph{(0).} For fixed $(I,J)$, $[LQL^\top]_{\mathrm{non\text{-}can}}(I,J)$ is
\[\sum_{\substack{(R_l,R_m,R_r):\text{side-inner-canonical, non-outer-canonical}\\\text{all three sizes}\le\tau}}(\omega/n)^{|V(R_l)|+|V(R_m)|+|V(R_r)|-|A|-|B|}q(R_m)\chi_{T_l\oplus T_m\oplus T_r},\tag{5.42}\]
where $(A,B)$ are the distinguished sets of $R_m$. Each triple has a unique separating factorization $(R'_l,R^*_m,R'_r)$.
\paragraph{First case: $|V(R^*_m)|\le\tau$.} The corresponding term in $(LQ'L^\top)_{\mathrm{can}}(I,J)$ is
\begin{multline*}(\omega/n)^{|V(R'_l)|+|V(R'_m)|+|V(R'_r)|-|S'_l|-|S'_r|}\\\cdot(\omega/n)^{z+|I(R^*_m)|}q(R'_m)\chi_{T'_l\oplus T^*_m\oplus T'_r},\tag{5.43}\end{multline*}
where $R'_m$ is the largest ribbon of $R^*_m$, $\chi_{T^*_m}$ means the character from $R'_m$, and $z\ge1$ is the original intersection size. Recall $T'_l\oplus T^*_m\oplus T'_r=T_l\oplus T_m\oplus T_r$ and
\begin{align*}|V(R_l)\cup V(R_m)\cup V(R_r)|&=|V(R'_l)|+|V(R^*_m)|+|V(R'_r)|-|S'_l|-|S'_r|\\&=|V(R_l)|+|V(R_m)|+|V(R_r)|-|A|-|B|-z.\end{align*}
Also $|V(R^*_m)|=|V(R'_m)|+|I(R^*_m)|$. Together the coefficient in (5.43) equals the one in (5.42) from $(R'_l,R^*_m,R'_r)$. Conversely, by $Q'$'s definition, (5.40) and Proposition5.24, each outer-canonical product corresponds uniquely to a side-inner-canonical triple of this case. Therefore the error collects precisely the next case.
\paragraph{Second case: $|V(R^*_m)|>\tau$.} We have
\begin{multline*}E'_{\mathrm{negl}}(I,J)=\sum_{\substack{(R_l,R_m,R_r):\text{side-inner-canonical, non-outer-canonical}\\\text{all three sizes}\le\tau,\ |V(R^*_m)|>\tau}}\\
(\omega/n)^{|V(R_l)|+|V(R_m)|+|V(R_r)|-|A|-|B|}\\
q(R_m)\chi_{T_l\oplus T_m\oplus T_r}.\tag{5.44}\end{multline*}
The obvious left/right distinguished-set condition $I,J$ is omitted.
\paragraph{(1).} For such a triple, putting $T=T_l\oplus T_m\oplus T_r$, the exponent is
\begin{align*}|V(R_l)|+|V(R_m)|+|V(R_r)|-|A|-|B|&=|V(R_l)\cup V(R_m)\cup V(R_r)|+z\\&=|V(T)\cup I\cup J|+|I(R^*_m)|+z.\end{align*}
Also $|I(R^*_m)|\le z+d/2$ from Lemma5.23.\footnote{In fact $|I(R^*_m)|\le z$ can be shown, but is not needed.} Fix $T$ and $a>\tau-|V(T)\cup I\cup J|$. To count triples contributing $(\omega/n)^{|V(T)\cup I\cup J|+a}\chi_T$ (temporarily ignore $q$), choose $I(R^*_m)$ of size $\le a/2+d/4$, since $a=|I(R^*_m)|+z\ge2|I(R^*_m)|-d/2$. Over fixed vertices the triple has $<3^{3\tau}2^{3\binom\tau2}$ choices. With $B_0=\max\{q(\cdot)\}$, its coefficient magnitude is less than
\begin{align*}&B_0(\omega/n)^{|V(T)\cup I\cup J|+a}n^{(a+d)/2}2^{2\tau^2}\\&=B_0(\omega/n^{1-2\epsilon})^{|V(T)\cup I\cup J|}n^{-2\epsilon|V(T)\cup I\cup J|}(\omega/\sqrt n)^an^{d/2}2^{2\tau^2}\\&\le B_0(n^{-1/2})^{|V(T)\cup I\cup J|}n^{-2\epsilon(|V(T)\cup I\cup J|+a)}n^{d/2}2^{2\tau^2}\\&\le B_0(n^{-1/2})^{|V(T)\cup I\cup J|}n^{-1.5\epsilon\tau}.\end{align*}
Use $\omega\le n^{1/2-4\epsilon}$, $|V(T)\cup I\cup J|+a>\tau$, $d<\epsilon\tau/10$, and $2^{2\tau}<n^{\epsilon/10}$. All appearing characters have $|V(T)|\le3\tau$. By Lemma4.2, for fixed $(I,J)$, w.p. $>1-n^{-10\log n}$,
\[|E'_{\mathrm{negl}}(I,J)|<\sum_{a=0}^{3\tau}B_0n^{-a/2}n^{-1.5\epsilon\tau}n^{a/2}n^{4\log\log n}2^{a^2}<n^{-1.4\epsilon\tau}.\]
Union-bound over $|\{(I,J)\}|<n^d$: w.p. $>1-n^{-9\log n}$, $\|E'_{\mathrm{negl}}\|<n^d n^{-1.4\epsilon\tau}<n^{-\epsilon\tau}$.
\paragraph{(2).} Fix $R_m$. By (5.40),
\begin{align*}q'(R_m)=\sum_{\substack{z,R^*_m:\text{largest ribbon }R_m}}(\omega/n)^{|I(R^*_m)|+z}\sum_{\substack{\mathcal P=(R_l,R,R_r):\text{side-inner-canonical}\\\mathcal P\to(R'_l,R^*_m,R'_r),\ z(\mathcal P)=z}}q(R).
\end{align*}
For fixed $R^*_m$ there are at most $8^{z\tau}<n^{\epsilon z}$ triples in the inner sum, since $(R'_l,R'_r)$ is fixed: after choosing each vertex's occurrence in the ribbons and $A,B\subseteq R^*_m$, only edges appearing more than once remain to assign. Each such edge has an endpoint in the fixed multiset $Z$ of size $\le z$. By Lemma5.23(2) and (5.41), the sum in absolute value is bounded by
\begin{align*}&n^{\epsilon z}(\omega/n)^{z+|I(R^*_m)|}|q(R)|\\&\le(\omega/n^{1-\epsilon})^{2(s'-s)+(p-p')+2|I(R^*_m)|}C(\omega/n^{1-\epsilon})^{s-p}\\&\le C(\omega/n)^{2|I(R^*_m)|}(\omega/n^{1-\epsilon})^{s'-p'+1/2}.\end{align*}
Here $(s,p)$ are the parameters of $R$ and $(s',p')$ those of $R_m$; the last step uses $s'-s\ge1/2$ from Lemma5.23(1). For fixed $i_0$ there are $<n^{i_0}$ choices of $R^*_m$ with $|I(R^*_m)|=i_0$, and $1\le z\le3\tau$. Hence
\begin{align*}|q'(R_m)|&\le3\tau\sum_{i_0=0}^{d/2}Cn^{i_0}(\omega/n)^{2i_0}(\omega/n^{1-\epsilon})^{s'-p'+1/2}\\&\le C(\omega/n^{1-\epsilon})^{s'-p'+1/3}.\end{align*}\end{proof}
Apply Lemma5.25 to the non-canonical term of (5.30), with $Q\leftarrow DQ_0D$: it becomes $[L(DQ_1D)L^\top]_{\mathrm{can}}+E'_{1;\mathrm{negl}}$. Repeat to obtain
\begin{align*}M'&=L[D(Q_0-Q_1+Q_2-\cdots\pm Q_d)D]L^\top\\&\quad-(E_{1;\mathrm{deg}}-\cdots\pm E_{1+d;\mathrm{deg}})+(E'_{1;\mathrm{negl}}+\cdots+E'_{d;\mathrm{negl}}).\tag{5.45}\end{align*}
\paragraph{Proposition 5.26 ([5],Claim6.15).} $Q_{d+1}=0$.
\begin{proof}Induct that every ribbon $R_m=(A,B;T_m)$ in $Q_k$ has $|A|+|B|\ge k$. The case $k=0$ is trivial. Every ribbon in $Q_{k+1}$ is the largest ribbon of a generalized middle in the separating factorization of a non-outer-canonical triple in $L(DQ_kD)L^\top$. The old middle has $|A|+|B|\ge k$, and Lemma5.23(1) gives $|A'|+|B'|\ge|A|+|B|+1\ge k+1$. For $k=d+1$ this is impossible with both distinguished sets in $\binom{[n]}{d/2}$.\end{proof}
\paragraph{Remark 5.27.} The final non-exact PSDness steps are omitted by the author, since similar arguments are given in the exact case, Section6. The selected target is the exact theorem; those complete Section6 steps follow below.

\section*{Author Section 6: PSDness of the exact pseudo-expectation}
Henceforth $M$ exclusively refers to the $d/2$-homogeneous minor of the moment matrix $\widetilde M$ in Definition3.13.
\paragraph{Theorem 6.1.} W.p. $>1-n^{-5\log n}$,
\[M(G)\succeq n^{-d-1}\operatorname{diag}(\widetilde{\mathrm{Cl}}(G))_{\binom{[n]}{d/2}\times\binom{[n]}{d/2}}.\]
\paragraph{Corollary 6.2.} W.p. $>1-n^{-5\log n}$, $\widetilde{\mathbb E}x_\emptyset>0$.
\begin{proof}By construction(3.13),
\[\widetilde{\mathbb E}x_\emptyset=\frac{\binom{\omega-d/2}{d-d/2}}{\binom{\omega-d}d\binom d{d/2}}\sum_{S:|S|=d/2}\widetilde{\mathbb E}x_S=\frac{\binom{\omega-d/2}{d-d/2}}{\binom{\omega-d}d\binom d{d/2}}\operatorname{Tr}(M),\]
positive w.h.p. by Theorem6.1.\end{proof}
Theorem1.4 is a quick corollary: the pseudo-expectation from Definition3.13 has PSD moment matrix by Theorem6.1 and Lemma4.1; Default holds by Corollary6.2 and rescaling as discussed before Remark3.9; Clique and Size Constraints hold by Lemma3.8. The degree-$d$ lower bound follows. The remainder proves Theorem6.1 by reducing it to Main Lemma6.9 and proving that lemma.
\subsection*{6.1 A Hadamard product and Euler transform}
In $M(I,J)$ the parameter $u=|I\cap J|$ prevents termwise factorization like the non-exact case. We express $M$ in a $\Sigma\Pi$-form where leaf matrices no longer depend on $u$, and factor those later.
\subsubsection*{6.1.1 Hadamard product}
Put $u=|I\cap J|$, $a=|V(T)\cup I\cup J|$. From(3.17),
\[M(I,J;T)=\sum_{c=0}^u\left[\frac{\omega^{u-c}}{\binom{\omega-d+u}u}\mathcal M_c(u,a)\right],\tag{6.1}\]
where
\[\mathcal M_c(u,a):=\binom{a-(d-u)}c\binom{n-a}{u-c}n^{-(u-c)}\frac{(a+u-c+8\tau^2)!}{(8\tau^2)!}(\omega/n)^a.\]
Thus
\[M=\sum_{c=0}^{d/2}m_c\circ M_c,\tag{6.2}\]
with matrices, for $|I|=|J|=d/2$,
\[m_c(I,J)=\frac{\omega^{u-c}}{\binom{\omega-d+u}u},\qquad u=|I\cap J|,\tag{6.3}\]
and
\[M_c(I,J)=\begin{cases}\sum_{T:|V(T)\cup I\cup J|\le\tau}\chi_T\mathcal M_c(|I\cap J|,|V(T)\cup I\cup J|),&|I\cap J|\ge c,\\0,&\text{o.w.}\end{cases}\tag{6.4}\]
\paragraph{Remark 6.3.} $m_c$ is supported on all $(I,J)$, while $M_c(I,J)=0$ when $|I\cap J|<c$. The intuition is that $M_c$'s are ``close'' for varying $c$, while $m_c$ qualitatively decreases in $c$. This would let us concentrate on $c=0$. The next lemma proves the second half; Lemma6.9 states the first more precisely.
\paragraph{Lemma 6.4.} For $c=0,\ldots,d/2$,
\[m_c=\omega^{-c}\sum_{k=0}^{d/2}b_kJ_k,\qquad b_k/k!\in[d/(2\omega),1+2dk/\omega],\]
where $J_k$ is (5.4). In particular,
\[m_0=\omega m_1=\cdots=\omega^{d/2}m_{d/2}\succ\frac1\omega\mathrm{Id}.\tag{6.5}\]
\begin{proof}By definition,
\[m_c=\omega^{-c}\sum_{l=0}^{d/2}\frac{\omega^l}{\binom{\omega-d+l}l}D_l.\]
By basis-change(5.6), $m_c=\omega^{-c}\sum_{k=0}^{d/2}J_kb_k$ with
\[b_k=k!\left(\sum_{l=0}^k(-1)^{k-l}\underbrace{\frac\omega{\omega-(d-l)}\cdots\frac\omega{\omega-(d-1)}\frac1{(k-l)!}}_{=:f_k(l),\text{ which is }1/k!\text{ if }l=0}\right).\]
For fixed $k$, $f_k(l)$ is increasing, so
\[\sum_{l=0}^k(-1)^{k-l}f_k(l)\ge f_k(k)-f_k(k-1)>\frac{d/2}\omega(1+(d/2)/\omega)^{k-1}\ge\frac d{2\omega}.\]
For $k=d/2$, $J_{d/2}=\mathrm{Id}$ gives(6.5).\end{proof}
\subsubsection*{6.1.2 Euler transform}
Fix $c$. For $(I,J;T)$ the coefficient $\mathcal M_c(u,a)$ is the partial function
\[\mathcal M_c(u,a)=\binom{a-(d-u)}c\binom{n-a}{u-c}n^{-(u-c)}\frac{(a+u-c+8\tau^2)!}{(8\tau^2)!}(\omega/n)^a.\tag{6.6}\]
\paragraph{Definition 6.5 (Extended $\mathcal M_c$).} It is defined on $u\ge c$, $u+a\ge d+c$ and extended to $\mathbb N^2$ by
\[\binom{n-a}{u-c}=0\text{ if }u<c,\tag{6.7}\]
and the usual conventions
\[\binom{-m}k=(-1)^k\binom{m+k-1}k\quad(0<m,\ 0\le k),\tag{6.8}\]
\[\binom mk=0\quad(0\le m<k).\tag{6.9}\]
In particular $\binom m0=1$ for all $m\in\mathbb Z$, and if $0\le a-(d-u)<c$, then $\mathcal M_c(u,a)=0$.
Seek a decomposition
\[M_c=\sum_{R\in\binom{[n]}{\le d/2}}M_c^R,\tag{6.10}\]
where $M_c^R$ is supported on row/column indices containing $R$. Explicitly,
\[M_c^R(I,J)=\begin{cases}\sum_{T:|V(T)\cup I\cup J|\le\tau}(\omega/n)^aY_c(|R|,a)\chi_T,&R\subseteq I,J,\\0,&\text{o.w.}\end{cases}\tag{6.11}\]
where $a=|V(T)\cup I\cup J|$ and $Y_c$ is to be chosen. Comparing each $(I,J;T)$, (6.10) is equivalent, for fixed $c,a$, to
\[\sum_{r=0}^u\binom urY_c(r,a)(\omega/n)^a=\mathcal M_c(u,a).\tag{6.12}\]
Thus take $Y_c(u,a)(\omega/n)^a$ to be the inverse Euler transform in $u$ of the extended $\mathcal M_c$.
\paragraph{Fact 6.6.} If sequences on $\mathbb N$ satisfy $x(m)=\sum_{l=0}^m\binom ml y(l)$, then $x$ is the Euler transform of $y$, whose inverse is $y(m)=\sum_{l=0}^m(-1)^{m-l}\binom ml x(l)$.\footnote{This fact is an application of the $\zeta$ matrix and its inverse.}
\paragraph{Definition 6.7 (Coefficients in $M_c^R$).}
\[Y_c(r,a)=\begin{cases}\sum_{l=c}^r(-1)^{r-l}\binom rl\binom{a+l-d}c\binom{n-a}{l-c}n^{-(l-c)}\dfrac{(a+l-c+8\tau^2)!}{(8\tau^2)!},&r\ge c,\\0,&\text{o.w.}\end{cases}\tag{6.13}\]
\paragraph{Lemma 6.8 (Hadamard-product decomposition).}
\[M=\sum_{c=0}^{d/2}m_c\circ\left(\sum_RM_c^R\right)=\sum_{R\in\binom{[n]}{\le d/2}}\underbrace{\left(\sum_{c=0}^{|R|}m_c\circ M_c^R\right)}_{=:M^R}.\tag{6.14--6.15}\]
Each $m_c$ is Lemma6.4, and:
\begin{enumerate}\item $M_c^R=0$ if $|R|<c$.
\item If $R\not\subseteq I\cap J$, $M_c^R(I,J)=0$.
\item If $|R|\ge c$ and $R\subseteq I\cap J$, then $M_c^R(I,J)=\sum_{T:|V(T)\cup I\cup J|\le\tau}M_c^R(I,J;T)\chi_T$, with
\begin{align*}M_c^R(I,J;T)&=(\omega/n)^a\sum_{l=c}^{|R|}(-1)^{|R|-l}\binom{|R|}l\binom{a+l-d}c\\&\quad\cdot\binom{n-a}{l-c}n^{-(l-c)}\frac{(a+l-c+8\tau^2)!}{(8\tau^2)!}\\&=(\omega/n)^aY_c(|R|,a).\tag{6.16}\end{align*}
\item For $0\le c\le r\le d/2$, $0\le a\le\tau$, $|Y_c(r,a)|<\tau^{5\tau}$.\end{enumerate}
\begin{proof}(1),(2),(3) are definition. To check $M_c=\sum_RM_c^R$, fix $(I,J;T)$ with $|I|=|J|=d/2$, $a=|V(T)\cup I\cup J|\le\tau$, and $u=|I\cap J|$. Note $a-(d-u)\ge0$ and
\[\sum_RM_c^R(I,J;T)=\sum_{R\subseteq I\cap J}M_c^R(I,J;T)=(\omega/n)^a\sum_{r=0}^{|I\cap J|}\binom{|I\cap J|}rY_c(r,a).\]
By Euler transform and(6.12) this equals extended $\mathcal M_c(u,a)$. It only remains to see $\mathcal M_c(u,a)=0$ if $u<c$ or $a-(d-u)<c$, by(6.7),(6.9). For(4),
\begin{align*}|Y_c(u,a)|&=\left|\sum_{l=c}^r(-1)^{r-l}\binom rl\binom{a+l-d}c\left[\binom{n-a}{l-c}n^{-(l-c)}\right]\frac{(a+l-c+8\tau^2)!}{(8\tau^2)!}\right|\\&<r\,2^r(2\tau)^r\cdot1\cdot(9\tau^2)^{2\tau}<\tau^{5\tau},\end{align*}
using $r\le d/2\ll\tau$ in our regime.\end{proof}
\paragraph{Lemma 6.9 (Main Lemma).} In(6.15), w.p. $>1-n^{-5\log n}$, for every $R\in\binom{[n]}{\le d/2}$, with $\mathcal P^R=\{I\in\binom{[n]}{d/2}:R\subseteq I\}$,
\[M_0^R\succeq n^{-d}\operatorname{diag}(\widetilde{\mathrm{Cl}})_{\mathcal P^R\times\mathcal P^R},\tag{6.17}\]
and
\[\pm\omega^{-c}M_c^R\preceq n^{-c/6}M_0^R\quad\forall0<c\le|R|.\tag{6.18}\]
\paragraph{Corollary 6.10 (Theorem6.1).} W.p. $>1-n^{-5\log n}$, $M(G)\succeq n^{-d-1}\operatorname{diag}(\widetilde{\mathrm{Cl}}(G))$ on $\binom{[n]}{d/2}\times\binom{[n]}{d/2}$.
\begin{proof}Suppose Lemma6.9 holds. For fixed $R$, $M^R=\sum_{c=0}^{|R|}m_c\circ M_c^R$. Hadamard product with a PSD matrix preserves PSDness (Schur product theorem), so
\begin{align*}\sum_{c=1}^{|R|}m_c\circ M_c^R&\preceq\sum_{c=1}^{|R|}m_c\circ(\omega^cn^{-c/6}M_0^R)\\&=\left(\sum_{c=1}^{|R|}n^{-c/6}m_0\right)\circ M_0^R\\&\preceq n^{-1/6}m_0\circ M_0^R.\end{align*}
Similarly the sum is $\succeq-n^{-1/6}m_0\circ M_c^R$. Thus
\[M^R\succeq(1-n^{-1/6})m_0\circ M_0^R\succeq n^{-d-1}\operatorname{diag}(\widetilde{\mathrm{Cl}})_{\mathcal P^R\times\mathcal P^R}\]
by Lemmas6.4/6.9(2). Apply(6.15):
\[M=M^\emptyset+\sum_{\emptyset\ne R\in\binom{[n]}{\le d/2}}M^R\succeq M^\emptyset\succeq n^{-d-1}\operatorname{diag}(\widetilde{\mathrm{Cl}})_{\binom{[n]}{d/2}\times\binom{[n]}{d/2}}.\tag{6.19}\]\end{proof}

\section*{Author Section 6.2: first-approximate factorization of $M_c^R$}
Use Section5.3's terminology and extend the recursive approximate factorization.
\paragraph{Definition 6.11.} Fix $R\in\binom{[n]}{\le d/2}$. For $i=0,\ldots,\tau$, define the left-$i$-factor $L_{R,i}$, of dimension $\binom{[n]}{d/2}\times\binom{[n]}{\le d/2}$, by
\[L_{R,i}(I,A)=\begin{cases}0,&R\not\subseteq I\cap A,\\\displaystyle\sum_{\substack{T:|V(T)\cup I\cup A|\le\tau,\ A=S_l(I,A;T)\\T\cap E(A)=\emptyset,\ (I,A;T)\text{ left-generated}\\e_{I,A}(T)=i}}(\omega/n)^i\chi_T,&\text{o.w.}\end{cases}\tag{6.20}\]
$(L_{R,j})^\top$ is the right-$j$-factor. $\widetilde L_R=(L_{R,0},\ldots,L_{R,\tau})$ is the left factor and $(\widetilde L_R)^\top$ the right factor. These do not depend on $c$.
\paragraph{Definition 6.12.} Let $D^\tau$ be the constant diagonal matrix
\[\operatorname{diag}((\omega/n)^{|A|/2})_{A:|A|\le d/2}\otimes\mathrm{Id}_{\{0,\ldots,\tau\}\times\{0,\ldots,\tau\}},\]
of dimension $(\binom{[n]}{\le d/2}\times(\tau+1))\times(\binom{[n]}{\le d/2}\times(\tau+1))$.
\paragraph{Definition 6.13 (Goal factorization).} Find a middle matrix $Q_c^R$ of that same dimension such that approximately
\[M_c^R\approx\widetilde L_R[D^\tau Q_c^R D^\tau](\widetilde L_R)^\top.\tag{6.21}\]
\paragraph{Remark 6.14.} Unlike Section5.3, further distinguish a parameter pair in $\{0,\ldots,\tau\}^2$: in(6.13), and any exact pseudo-expectation generated by Section3.2, $a=|V(T)\cup I\cup J|$ appears nestedly. Previously $(\omega/n)^a=(\omega/n)^{e(R_l)+|V(R_m)|+e(R_r)}$ naturally factors left/middle/right. Here terms like $\binom{a+l-d}c\binom{n-a}{l-c}$ are not log-additive in $a$. The chosen generating function makes $\mathbb E[Q_{0,0}^R]$ positive in this harder situation, as Proposition6.28/Corollary6.30 will show. Separate the main factor $(\omega/n)^a$ as before, and leave $Y_c(r,a)$ to the middle matrix indexed by $(e_l,e_r)$.
\paragraph{Definition 6.15 (First-approximate factorization).} $Q_{c,0}^R$ is a $(\tau+1)\times(\tau+1)$ block matrix, each block of dimension $\binom{[n]}{\le d/2}\times\binom{[n]}{\le d/2}$, zero outside the principal minor $\mathcal S^R\times\mathcal S^R$ with
\[\mathcal S^R:=\{(A,i)\in\binom{[n]}{\le d/2}\times\{0,\ldots,\tau\}:A\supseteq R,\ |A|+i\ge d/2\}.\tag{6.22}\]
Inside this minor,
\begin{align*}Q_{c,0}^R((A,i),(B,j))=\sum_{\substack{T_m:|V(T_m)\cup A\cup B|\le\tau\\A,B\in\mathrm{mSep}_{A,B}(T_m)}}&(\omega/n)^{|V(T_m)\cup A\cup B|-(|A|+|B|)/2}\\&\cdot Y_c(|R|,|V(T_m)\cup A\cup B|+(i+j))\chi_{T_m}.\tag{6.23}\end{align*}
Define $\widetilde L_R[D^\tau Q_{c,0}^R D^\tau](\widetilde L_R)^\top$ as the first approximate factorization of $M_c^R$.
\paragraph{Remark 6.16.} (1) $\mathcal S^R$ is independent of $c$; $|A|+i\ge d/2$ is natural since $i$ is intended as $|V(T')\setminus A|\ge|I|-|A|$ for a left ribbon. Otherwise that column in $\widetilde L_R$ is zero. Similarly for $j$. (2) $Q_{c,0}^R$ is supported where $|A|=|B|$. (3) For an outer-canonical product in the above matrix product, Remark5.12 gives $|V(T_m)\cup A\cup B|+(i+j)=|V(T)\cup I\cup J|$ for the canonical decomposition with $e(R_l)=i$, $e(R_r)=j$.
\paragraph{Lemma 6.17.} Fix $R,c\le|R|$. For each $(I,J;T)$ with $|V(T)\cup I\cup J|\le\tau$ and $R\subseteq I\cap J$, there is exactly one outer-canonical product in
\[\widetilde L_R[D^\tau Q_{c,0}^R D^\tau](\widetilde L_R)^\top\tag{6.24}\]
corresponding to its canonical decomposition, giving the term $M_c^R(I,J;T)\chi_T$.
\begin{proof}Every triple in(6.24) is inner-canonical by the definitions of $\widetilde L_R,Q_{c,0}^R$, hence outer-canonical triples 1-1 correspond to $(I,J;T)$ by canonical decomposition. Fix such a decomposition. $(I,A;T')$ appears once in $\widetilde L_R(I,A)$ in block $L_{R,e_l}$, $e_l=e_{I,A}(T')$; symmetrically for $(J,B;T'')$, $e_r=e_{J,B}(T'')$. Its unique outer-canonical product coefficient is
\[L_{R,e_l}(I,A;T')(\omega/n)^{|A|/2}Q_{c,0}^R(A,B;T_m)(\omega/n)^{|B|/2}L_{R,e_r}(J,B;T'').\tag{6.25}\]
By(6.20),(6.23), if $a=|V(T)\cup I\cup J|\le\tau$, this equals $(\omega/n)^aY_c(|R|,a)=M_c^R(I,J;T)$ by(6.13),(6.16), since canonicality gives $a=e_l+|V(T_m)\cup A\cup B|+e_r$.\end{proof}
\paragraph{Definition 6.18 (First error-matrices).} Let $E_{c,1;\mathrm{negl}}$ sum all outer-canonical products in(6.24) exceeding degree, namely $|V(T)\cup I\cup J|>\tau$. Let the non-canonical part sum all non-outer-canonical products. Lemma6.17 says $M_c^R=[\widetilde L_R(D^\tau Q_{c,0}^R D^\tau)(\widetilde L_R)^\top]_{\mathrm{can}}$, equivalently
\begin{align*}M_c^R&=\widetilde L_R(D^\tau Q_{c,0}^R D^\tau)(\widetilde L_R)^\top\\&\quad-[\widetilde L_R(D^\tau Q_{c,0}^R D^\tau)(\widetilde L_R)^\top]_{\mathrm{non\text{-}can}}-E^R_{c,1;\mathrm{deg}}.\tag{6.26}\end{align*}
The crucial fact is that the main error approximately factorizes through the same left/right factors; the next subsection extends recursive factorization to do this.
\section*{Author Section 6.3: exact recursive factorization statement}
\paragraph{Lemma 6.19 (Recursive approximate factorization; exact case).} For any fixed $R\in\binom{[n]}{\le d/2}$ and $0\le c\le|R|$,
\[M_c^R=\widetilde L_R[D^\tau(Q_{c,0}^R-Q_{c,1}^R+\cdots\pm Q_{c,d}^R)D^\tau](\widetilde L_R)^\top+E_c^R,\tag{6.27}\]
where:
\begin{enumerate}\item All $Q_{c,k}^R$ are supported on $\mathcal S^R\times\mathcal S^R$, with $\mathcal S^R$ from(6.22).
\item $Q_{c,0}^R$ is Definition6.15.
\item For $1<k\le d/2$, within that minor the $(i,j)$ block is
\[Q_{c,k}^R((A,i),(B,j))=\sum_{T_m:|V(T_m)\cup A\cup B|\le\tau}q_{c,k}^R(R_m,i,j)\chi_{T_m},\tag{6.28}\]
$R_m=(A,B;T_m)$, the coefficients symmetric w.r.t. shapes, and
\[\forall(i,j)\quad|q_{c,k}^R(R_m,i,j)|\le\tau^{5\tau}(\omega/n^{1-\epsilon})^{s-p+k/3},\tag{6.29}\]
where $s=(|A|+|B|)/2$ and $p$ is the maximum number of vertex-disjoint paths from $A$ to $B$ in $R_m$.
\item For any $G$, $E_c^R(G)$ is supported on clique rows/columns containing $R$, and w.p. $>1-n^{-9\log n}$,
\[\|E_c^R\|<n^{-\epsilon\tau/2}.\tag{6.30}\]
\end{enumerate}
The proof proceeds via one-round Lemma6.20, followed by its iteration; the complete author proof is continued in the next fragment.

\section*{Complete author exact recursive-factorization proof}
Again fix $R\in\binom{[n]}{\le d/2}$, $c\le|R|$, and write $n_1:=\binom{[n]}{d/2}\times(\tau+1)$ for convenience.
\paragraph{Lemma 6.20 (One round of factorization; exact case).} Let $\widetilde L_R$ be Definition6.11 and $Q^R$ any $n_1\times n_1$ matrix supported on $\mathcal S^R\times\mathcal S^R$, with
\[Q^R((A,i),(B,j))=\sum_{T_m:|V(T_m)\cup A\cup B|\le\tau}(\omega/n)^{|V(R_m)|}q(R_m,i,j)\chi_{T_m},\tag{6.31}\]
where $R_m=(A,B;T_m)$, $q(\cdot,i,j)$ symmetric w.r.t. shapes for fixed $(i,j)$. Define $Q'$ and $E'_{\mathrm{negl}}$ so that
\[[\widetilde L_RQ(\widetilde L_R)^\top]_{\mathrm{non\text{-}can}}=[\widetilde L_RQ'(\widetilde L_R)^\top]_{\mathrm{can}}+E'_{\mathrm{negl}}.\tag{6.32}\]
Let $Q'$ be supported on $\mathcal S^R\times\mathcal S^R$ and
\[Q'((A,i),(B,j))=\sum_{T_m:|V(T_m)\cup A\cup B|\le\tau}(\omega/n)^{|V(R_m)|}q'(R_m,i,j)\chi_{T_m}.\tag{6.33}\]
For fixed $R_m=(A,B;T_m)$, $(i,j)$, put $t=|V(R_m)|\le\tau$, $s=(|A|+|B|)/2$. For every generalized $R^*_m$ with largest ribbon $R_m$ and size $\le\tau$, fix a pair $(R'_l,R'_r)$ such that $(R'_l,R^*_m,R'_r)$ is a separating factorization of some triple, $|V(R'_l)|,|V(R'_r)|\le\tau$ and
\[(e(R'_l),e(R'_r))=(i,j).\tag{6.34}\]
Exclude $R^*_m$ if no choice exists. Set
\begin{align*}q'(R_m,i,j)&=\sum_{\substack{R^*_m:\text{generalized on }(A,B)\\|V(R^*_m)|\le\tau,\ \text{largest ribbon }R_m}}(\omega/n)^{|I(R^*_m)|}q''(R^*_m,i,j),\\q''(R^*_m,i,j)&=\sum_{\substack{z,i_1,j_1:\ 1\le z\le d/2}}\sum_{\substack{\mathcal P=(R_l,R,R_r):\text{side-inner-canonical}\\\mathcal P\to(R'_l,R^*_m,R'_r),\ z(\mathcal P)=z\\e(R_l)=i_1,\ e(R_r)=j_1}}(\omega/n)^zq(R,i_1,j_1).\tag{6.35}\end{align*}
The full Proposition5.24 makes this independent of the fixed pair, so $q'$ is symmetric w.r.t. shapes. Define the error by(6.32). Conclusions: (1) w.p. $>1-n^{-9\log n}$, $\|E'_{\mathrm{negl}}\|\le\max\{q(\cdot)\}n^{-\epsilon\tau}$. (2) If
\[\forall R_m,i,j\quad|q(R_m,i,j)|\le C(\omega/n^{1-\epsilon})^{s-p},\tag{6.36}\]
where $p$ is maximum vertex-disjoint $A$--$B$ path count, then $|q'(R_m)|\le C(\omega/n^{1-\epsilon})^{s-p+1/3}$ for all $R_m,i,j$.
\paragraph{Proof of Lemma 6.20.} The proof is almost the same as Lemma5.25; the author gives all differences as follows. The support condition does not affect anything, since $\widetilde L_R$ is automatically zero on columns/rows not in $\mathcal S^R$. Expand the canonical product with $Q'$ to compare termwise to the non-canonical product with $Q$, by Proposition5.24. When $(i,j)$ and $R^*_m$ are fixed, the sizes of any pair satisfying(6.34) are fixed, so that proposition applies. The comparison of orders in $\omega/n$ is exactly step(0) of Lemma5.25. Thus the error collects all non-canonical terms whose separating middle exceeds size $\tau$:
\begin{multline*}E'_{\mathrm{negl}}(I,J)=\sum_{i,j}\sum_{\substack{(R_l,R_m,R_r):\text{side-inner-canonical, non-outer-canonical}\\\text{all three sizes}\le\tau,\ |V(R^*_m)|>\tau\\(e(R_l),e(R_r))=(i,j)}}\\
(\omega/n)^{|V(R_l)|+|V(R_m)|+|V(R_r)|-|A|-|B|}\\\cdot q(R_m,i,j)\chi_T,\\
\qquad T=T_l\oplus T_m\oplus T_r.\tag{6.37}\end{multline*}
Left/right distinguished sets $I,J$ are implicit. Numerical conclusions(1),(2) follow from the estimates after(5.44). For(1), that estimate is loose enough that even an extra $(1+\tau)^2$ factor from union-bound over blocks remains below $n^{-\epsilon\tau}$.\qed
\paragraph{Proof of Lemma 6.19.} Apply Lemma6.20 to $[\widetilde L_R(D^\tau Q_{c,i}^R D^\tau)(\widetilde L_R)^\top]_{\mathrm{non\text{-}can}}$ for $i=0$, obtaining $Q_{c,1}^R,E'_{1;\mathrm{negl}}$ (hide $R,c$ on this negligible matrix). Repeat for $i=1$ obtaining $E_{c,1;\mathrm{deg}},Q_{c,2}^R,E'_{2;\mathrm{negl}}$. Continuing gives
\begin{align*}M_c^R&=\widetilde L_R[D^\tau(Q_{c,0}^R-Q_{c,1}^R+\cdots\pm Q_{c,d}^R)D^\tau](\widetilde L_R)^\top\\&\quad-(E_{c,1;\mathrm{deg}}^R-E_{c,2;\mathrm{deg}}^R+\cdots\pm E_{c,d;\mathrm{deg}}^R)\\&\quad+(E'_{1;\mathrm{negl}}+\cdots+E'_{d;\mathrm{negl}}).\tag{6.38}\end{align*}
Here $Q_{c,d+1}^R=0$ by the same Proposition5.26.
\begin{enumerate}\item All $Q_{c,k}^R$ have the required support by the definition of each round.
\item By definition.
\item Each coefficient family remains symmetric w.r.t. shapes by Lemma6.20. From(6.23) and Lemma6.8(4), $|q_{c,0}^R(R_m)|=|Y_c(|R|,|R_m|)|\le\tau^{5\tau}$. Every ribbon appearing in $Q_{c,0}^R$ has $|A|=|B|$ vertex-disjoint $A$--$B$ paths, so $s=p$. Equivalently,
\[|q_{c,0}^R(R_m)|\le(\omega/n^{1-\epsilon})^{s-p}\tau^{5\tau}.\tag{6.39}\]
Use Lemma6.20(2): its $q$ corresponds to $q_{c,k}^R$, because its matrix $Q$ is $D^\tau Q_{c,k}^R D^\tau$, whose $(\omega/n)^{|V(R_m)|}q$ equals $(\omega/n)^{|V(R_m)|-s}(\omega/n)^s q_{c,k}^R$. Thus recursively
\[|q_{c,k}^R(R_m,i,j)|\le\tau^{5\tau}(\omega/n^{1-\epsilon})^{s-p+k/3}.\]
\item For every $G$, both $M_c^R$ and its factored expression are supported on clique rows/columns containing $R$, so their difference $E_c^R$ is too. From(6.38) it is the sum of alternating degree errors and negligible errors. By Lemma6.20(2), inductively all $|q_{c,k}^R|<\tau^{5\tau}$. By Lemma6.20(1), each $E'_{k;\mathrm{negl}}$ satisfies w.p. $>1-n^{-9\log n}$,
\[\|E'_{k;\mathrm{negl}}\|<\tau^{5\tau}n^{-\epsilon\tau}<n^{-0.9\epsilon\tau}.\]
\end{enumerate}
For the degree error, Definition5.15 says $E_{c,k;\mathrm{deg}}^R(I,J)$ sums all outer-canonical products exceeding degree in $\widetilde L_R(D^\tau Q_{c,i-1}^R D^\tau)(\widetilde L_R)^\top$. Thus
\begin{multline*}E_{c,k;\mathrm{deg}}^R(I,J)=\sum_{\substack{(R_l,R_m,R_r):\text{semi-inner-canonical, outer-canonical}\\|V(T)\cup I\cup J|>\tau}}(\omega/n)^{|V(T)\cup I\cup J|}\\\cdot q_{c,k-1}^R(R_m,e(R_l),e(R_r))\chi_T.\end{multline*}
As usual $T=T_l\oplus T_m\oplus T_r$, side sets are $I,J$, and $s$ is the average size of the middle distinguished sets. Outer/semi-inner canonicality gives $|V(T)\cup I\cup J|=e_l+e_r+|V(R_m)|$. Each fixed $(I,J;T)$ can come from at most $3^{3\tau}$ triples, since their vertex-set union is exactly $V(T)\cup I\cup J$. Since $3\tau\ge|V(T)\cup I\cup J|>\tau$ and w.h.p. $|q_{c,k-1}^R|<\tau^{5\tau}$, Lemma4.2 gives w.p. $>1-n^{-10\log n}$,
\[|E_{c,k;\mathrm{deg}}^R(I,J)|<\tau^{6\tau}\sum_{c=0}^{3\tau}(\omega/n)^{\max\{\tau,c\}}n^{c/2}2^{c^2}n^{4\log\log n}<n^{-2\epsilon\tau}.\]
Union-bound over $(I,J)$ gives $\|E_{c,k;\mathrm{deg}}^R\|<n^{-d/4}n^{-2\epsilon\tau}<n^{-\epsilon\tau}$ w.p. $>1-n^{-9.5\log n}$. Sum the two errors and union-bound over $k$: w.p. $>1-n^{-9\log n}$, $\|E_c^R\|<n^{-\epsilon\tau/2}$.\qed

\section*{Author Sections 6.4--6.5: positivity of middle matrices}
Write $Q_c^R=Q_{c,0}^R-Q_{c,1}^R+\cdots\pm Q_{c,d}^R$ in Lemma6.19's approximate decomposition. The following is the key lemma, proved through the next subsections.
\paragraph{Lemma 6.21.} W.p. $>1-n^{-8\log n}$: (1) for all $R\in\binom{[n]}{\le d/2}$,
\[Q_{0,0}^R-Q_{0,1}^R+\cdots\pm Q_{0,d/2}^R\succeq\tau^{-7\tau}\operatorname{diag}(\widetilde{\mathrm{Cl}})_{\mathcal S^R\times\mathcal S^R};\]
(2) for all $R$, $0<c\le|R|$,
\[\pm\omega^{-c}(Q_{c,0}^R-Q_{c,1}^R+\cdots\pm Q_{c,d/2}^R)\preceq n^{-c/4}\operatorname{diag}(\widetilde{\mathrm{Cl}})_{\mathcal S^R\times\mathcal S^R}.\]
The proof has three ingredients: Corollary6.36 shows positivity of $Q_{0,0}^R$ via a structural/pseudorandom decomposition; Lemmas6.37/6.38 bound the rest. The last real technical challenge is $\mathbb E[Q_{0,0}^R]$, ultimately the choice of $F$ in Definition3.11.
\paragraph{Definition 6.22.} The root diagonal-clique matrix on $\binom{[n]}{\le d/2}$ is
\[D_{\mathrm{Cl}}(A,B)=\begin{cases}0,&A\ne B,\\2^{-\binom{|A|}2/2}\widetilde{\mathrm{Cl}}_A=2^{-\binom{|A|}2/2}\sum_{T\subseteq E[A]}\chi_T,&A=B.\end{cases}\tag{6.40}\]
Thus $D_{\mathrm{Cl}}^2(A,A)=\widetilde{\mathrm{Cl}}(A)$. Define $D_{\mathrm{Cl}}^\tau:=D_{\mathrm{Cl}}\otimes\mathrm{Id}_{\{0,\ldots,\tau\}\times\{0,\ldots,\tau\}}$ (6.41), also diagonal.
\paragraph{Definition 6.23.} The structural-pseudorandom decomposition is
\[Q_{0,0}^R=D_{\mathrm{Cl}}^\tau\mathbb E[Q_{0,0}^R]D_{\mathrm{Cl}}^\tau+\left(Q_{0,0}^R-D_{\mathrm{Cl}}^\tau\mathbb E[Q_{0,0}^R]D_{\mathrm{Cl}}^\tau\right).\tag{6.42}\]
The first summand is structural and the second pseudorandom.
\paragraph{Proposition 6.24 (Expression of $\mathbb E[Q_{c,0}^R]$).} Fix $R$, $0\le c\le|R|$, $r=|R|$.
\begin{enumerate}\item The expectation is supported on blockwise partial diagonals $\{((A,i),(A,j))\in\mathcal S^R\times\mathcal S^R\}$, requiring $R\subseteq A$ and $|A|+\min\{i,j\}\ge d/2$.
\item On these positions,
\begin{align*}\mathbb E[Q_{c,0}^R]((A,i),(A,j))&=\sum_{l=c}^r(-1)^{r-l}\frac{\binom rl}{(l-c)!}\binom{|A|+i+j+l-d}c\\&\quad\cdot\frac{(|A|+8\tau^2+(l-c)+(i+j))!}{(8\tau^2)!}+O(\tau^{1.5\tau}/n).\tag{6.43}\end{align*}
For $c=0$ this is
\[\sum_{l=0}^r(-1)^{r-l}\frac{\binom rl}{l!}\frac{(|A|+8\tau^2+l+(i+j))!}{(8\tau^2)!}+O(\tau^{1.5\tau}/n).\tag{6.44}\]
\item If $1_{A,A}$ is the matrix with one 1 at $(A,A)$, then
\[\mathbb E[Q_{0,0}^R]=\sum_{\substack{A\in\binom{[n]}{\le d/2}\\A\supseteq R}}1_{A,A}\otimes\left[\sum_{l=0}^r(-1)^{r-l}\frac{\binom rl}{l!}P_{|A|+l}+E_A^R\right].\tag{6.45}\]
$P_{|A|+l}$ and $E_A^R$ are $(\tau+1)\times(\tau+1)$ matrices supported on $\{i:d/2-|A|\le i\le\tau\}^2$,
\[\|E_A^R\|<\tau^{2\tau}/n,\tag{6.46}\]
and
\[P_{|A|+l}(i,j)=\frac{(|A|+l+8\tau^2+(i+j))!}{(8\tau^2)!}.\tag{6.47}\]
\end{enumerate}
\begin{proof}(1) Constant terms in(6.23) have $T_m=\emptyset$, nonzero only when $A=B$. (2) For those positions the expectation is $Y_c(r,a)$, $a=|A|+i+j$, namely
\[\sum_{l=c}^r(-1)^{r-l}\binom rl\binom{a+l-d}c\binom{n-a}{l-c}n^{-(l-c)}\frac{(a+l-c+8\tau^2)!}{(8\tau^2)!}.\tag{6.48}\]
Use
\[\binom{n-a}{l-c}n^{-(l-c)}=\frac1{(l-c)!}\frac{(n-a)\cdots(n-a-(l-c)+1)}{n^{l-c}}=\frac1{(l-c)!}(1-O(d^2/n)),\]
and the bound
\[\binom rl\binom{a+l-d}c\binom{n-a}{l-c}n^{-(l-c)}\frac{(a+l-c+8\tau^2)!}{(8\tau^2)!}<(4d)^d(9\tau^2)^d<\tau^\tau\]
to obtain(6.43). For $c=0$, the author writes that $\binom{a+l-d}0=0$ regardless of $a+l-d$ in(6.48), and the same analysis gives(6.44). (3) Each error matrix has dimension $(\tau+1)^2$ and entries absolutely $<\tau^{1.5\tau}/n$ by(2).\end{proof}
\paragraph{Remark 6.25.} The specialty of $c=0$ in(6.44) is that $\binom{|A|+l-d}0$ is always 1. For $c>0$, $\binom{|A|+l-d}c$ can be zero or negative depending on the order between $0,c,|A|+l-d$, so its expected matrix may not be PSD.
\paragraph{Definition 6.26.} For $m,t\in\mathbb N$, the factorial Hankel matrix is $H_{m,t}(i,j)=(i+j+t)!$ for $0\le i,j\le m$ (6.49).
\paragraph{Proposition 6.27 (Almost common decomposition).} (1)
\[H_{m,t}=L_m[N_{m,t}D_{m,t}(N_{m,t})^\top]L_m^\top,\]
where $L_m,D_{m,t}$ are diagonal and $N_{m,t}$ lower-triangular,
\[L_m(i,i)=i!,\quad D_{m,t}(i,i)=\prod_{t'=1}^t(i+t'),\quad N_{m,t}(i,j)=\binom{i+t}{i-j}.\]
$L_m$ is independent of $t$ and $H_{m,t}$ positive. (2) Let $J_m$ be the $(m+1)\times(m+1)$ lower-triangular Jordan block, equal to 1 if $i=j$ or $i=j+1$ and 0 otherwise. Then $N_{m,t+1}=N_{m,t}J_m$ (6.50).
\begin{proof}This follows from a direct inspection.\end{proof}
\paragraph{Proposition 6.28.} If $t+1>8\max\{r^2,m\}$ (6.51), then $H_{m,t+1}\succeq2r^2H_{m,t}$.
\begin{proof}By Proposition6.27 it suffices to show $J_mD_{m,t+1}J_m^\top\succeq2r^2D_{m,t}$. For $x^\top=(x_0,\ldots,x_m)$ compare
\[q_{t+1}(x):=(x^\top J_m)D_{m,t+1}(J_m^\top x),\qquad q_t(x):=2r^2x^\top D_{m,t}x.\tag{6.52}\]
Put $\alpha(y)=2r^2\prod_{t'=1}^t(y+t')$, $\beta(y)=\prod_{t'=1}^{t+1}(y+t')$. Then $q_{t+1}(x)=\sum_{i=0}^m\beta(i)(x_i+x_{i+1})^2$, $x_{m+1}:=0$, while $q_t(x)=\sum_{i=0}^m\alpha(i)x_i^2$. Completing squares gives
\begin{multline*}q_{t+1}(x)=\sum_{i=0}^m\left[\alpha(i)x_i^2+(\beta(i)-\alpha(i))\left(x_i+\frac{\beta(i)}{\beta(i)-\alpha(i)}x_{i+1}\right)^2\right.\\\left.-\frac{\beta(i)^2}{\beta(i)-\alpha(i)}x_{i+1}^2\right].\end{multline*}
Define $b_0=1-\alpha(0)/\beta(0)$ and
\[b_i=1-\frac{\alpha(i)}{\beta(i)}-\frac{\beta(i-1)}{\beta(i)}\frac1{b_{i-1}}.\tag{6.53}\]
Then
\[q_{t+1}(x)=\underbrace{\sum_{i=0}^m\alpha(i)x_i^2}_{q_t(x)}+\sum_{i=0}^m\beta(i)b_i\left(x_i+\frac1{b_i}x_{i+1}\right)^2.\tag{6.54}\]
\paragraph{Claim 6.29.} Each $b_i>1/2$ for $i\le m$.
\paragraph{Proof of claim.} $b_0=1-2r^2/(t+1)$ and, for $i\ge1$,
\[b_i=1-\frac{2r^2}{t+1+i}-\frac i{t+1+i}\frac1{b_{i-1}}.\tag{6.55}\]
Induct: $b_0>1/2$ by(6.51), and
\[b_i\ge1-\frac{2r^2}{t+1}-\frac m{t+1}\cdot2>1/2\]
by(6.51) and the inductive hypothesis. This proves the claim. Since $b_i$ are positive, (6.54) gives $q_{t+1}\ge q_t$, proving the proposition.\end{proof}
Apply this to(6.47): on its support,
\[P_{|A|+l}=\frac1{(8\tau^2)!}H_{\tau-(d/2-|A|),\ d-|A|+8\tau^2+l}.\]
\paragraph{Corollary 6.30 (Positiveness of $\mathbb E[Q_{0,0}^R]$).} In(6.45), regarding matrices on their support,
\[\sum_{l=0}^r(-1)^{r-l}\frac{\binom rl}{l!}P_{|A|+l}+E_A^R\succ\operatorname{diag}(\tau^{-6\tau})_{0\le i\le\tau-(d/2-|A|)}.\tag{6.56}\]
In particular,
\[\mathbb E[Q_{0,0}^R]\succ\sum_{A\supseteq R}1_{A,A}\otimes\operatorname{diag}(\tau^{-6\tau})_{d/2-|A|\le i\le\tau}=\operatorname{diag}(\tau^{-6\tau})_{\mathcal S^R\times\mathcal S^R}.\tag{6.57}\]
\begin{proof}The last part follows by checking support and that tensoring with a nonzero PSD matrix preserves $\succ$. For(6.56), fix $A$ and put
\[\tau_0=\tau-(d/2-|A|),\qquad t_0=d-|A|+8\tau^2.\tag{6.58}\]
Then
\[\sum_{l=0}^r(-1)^{r-l}\frac{\binom rl}{l!}P_{|A|+l}=\frac1{(8\tau^2)!}(X_r+X_{r-2}+\cdots),\tag{6.59}\]
where for $0\le v\le\lfloor r/2\rfloor$,
\[X_{r-2v}=\frac{\binom r{r-2v}}{(r-2v)!}\left(H_{\tau_0,t_0+r-2v}-\frac{(r-2v)^2}{2v+1}H_{\tau_0,t_0+r-2v-1}\right),\quad H_{\tau_0,-1}:=0.\]
The second coefficient is $\le r^2$. As $t_0>8\max\{r^2,\tau_0\}$, Proposition6.28 gives
\[X_{r-2v}\succeq\frac{\binom r{r-2v}}{(r-2v)!}\max\{\tfrac12H_{\tau_0,t_0+r-2v},\ r^2H_{\tau_0,t_0+r-2v-1}\}.\]
In particular,
\begin{align*}\sum_{l=0}^r(-1)^{r-l}\frac{\binom rl}{l!}P_{|A|+l}&\succeq\frac1{(8\tau^2)!}H_{\tau_0,t_0}\\&=L[N_{t_0}\tfrac{D_{t_0}}{(8\tau^2)!}(N_{t_0})^\top]L^\top,\tag{6.60}\end{align*}
omitting index $\tau_0$ on the RHS. Claim6.31 below implies the RHS is $\succ L\operatorname{diag}(\tau^{-5\tau})L\succeq\operatorname{diag}(\tau^{-5\tau})$. By Proposition6.24(3), $\|E_A^R\|<\tau^{2\tau}/n<\tau^{-6\tau}$ in the parameter regime. Thus the LHS of(6.56) is $\succeq\operatorname{diag}(\tau^{-5\tau}-\tau^{-6\tau})\succeq$ its RHS.\end{proof}
\paragraph{Claim 6.31.} In Corollary6.30's notation,
\[N_{t_0}^{-1}(i,j)=(-1)^{i-j}\binom{i+t_0}{i-j},\quad0\le i,j\le\tau_0,\tag{6.61}\]
and
\[N_{t_0}\frac{D_{t_0}}{(8\tau^2)!}(N_{t_0})^\top\succ\operatorname{diag}(\tau^{-5\tau})_{0\le i\le\tau_0}.\tag{6.62}\]
\begin{proof}For(6.61), multiplying by $N_{t_0}$ gives at $(i,j)$
\begin{align*}\sum_{j\le k\le i}(-1)^{i-k}\binom{i+t_0}{i-k}\binom{k+t_0}{k-j}&=\sum_{k'=0}^{i'}(-1)^{i'-k'}\binom{i'+j+t_0}{i'-k'}\binom{k'+j+t_0}{k'},\end{align*}
where $i'=i-j$, $k'=k-j$. To see identity, let $D_m[(1+x)^a]$ be the coefficient of $x^m$ for $m\ge0,a\in\mathbb Z$. The RHS is
\begin{align*}&(-1)^{i'}\sum_{k'=0}^{i'}D_{i'-k'}[(1+x)^{i'+j+t_0}]D_{k'}[(1+x)^{-(t_0+j+1)}]\\&=(-1)^{i'}D_{i'}[(1+x)^{i'+j+t_0-(t_0+j+1)}]\\&=(-1)^{i'}D_{i'}[(1+x)^{i'-1}]=\mathbf1_{i'=0}.
\end{align*}
For(6.62), equivalently
\[\frac{D_{t_0}}{(8\tau^2)!}\succ N_{t_0}^{-1}\tau^{-5\tau}(N_{t_0}^{-1})^\top.\tag{6.63}\]
Set $a_0=\tau^{-5\tau}$. The RHS quadratic form is
\[x^\top N_{t_0}^{-1}a_0(N_{t_0}^{-1})^\top x=a_0\sum_{j=0}^{\tau_0}y_j^2,\tag{6.64}\]
with $y_j=(x^\top N_{t_0}^{-1})_j=\sum_{i=j}^{\tau_0}(-1)^{i-j}\binom{i+t_0}{i-j}x_i$. By Cauchy--Schwarz, $y_j^2\le\tau_0\sum_{i=j}^{\tau_0}\binom{i+t_0}{i-j}^2x_i^2$. Hence
\begin{align*}a_0\sum_jy_j^2&\le a_0\sum_{i=0}^{\tau_0}x_i^2\left(\tau_0\sum_{j=0}^i\binom{i+t_0}{i-j}^2\right)\\&<\sum_{i=0}^{\tau_0}\tau^{-5\tau}(9\tau^2)^{2i+2}x_i^2.\end{align*}
But the LHS diagonal of(6.63) satisfies
\[\frac{D_{t_0}(i,i)}{(8\tau^2)!}\ge(8\tau^2)^{-(d/2-|A|)}>\tau^{-2d}>\tau^{-5\tau}(9\tau^2)^{2i+2},\]
using $i\le\tau_0<\tau$ and $d\ll\tau$. Thus(6.63) holds.\end{proof}
\paragraph{Corollary 6.32 (Positive structural part).}
\[D_{\mathrm{Cl}}^\tau\mathbb E[Q_{0,0}^R]D_{\mathrm{Cl}}^\tau\succeq\tau^{-6\tau}\operatorname{diag}(\widetilde{\mathrm{Cl}})_{\mathcal S^R\times\mathcal S^R}.\]
\begin{proof}This follows from Corollary6.30 and $D_{\mathrm{Cl}}^2(A,A)=\widetilde{\mathrm{Cl}}(A)$ in Definition6.22.\end{proof}

\section*{Author Section 6.6: bounds on the remaining matrices}
We bound the pseudorandom part of $Q_{0,0}^R$, all $Q_{0,k}^R$ for $k>0$, and all $\omega^{-c}Q_{c,k}^R$ for $c>0$, by Lemmas6.34,6.37,6.38. This proves Lemma6.21. Care is required on the block structure and supports.
\paragraph{Definition 6.33 (0-1 diagonal-clique matrix).} $D'$ is the 0-1 version of $D_{\mathrm{Cl}}^\tau$, diagonal with $D'((A,i),(A,i))=\mathrm{Cl}_A$ for all $A\in\binom{[n]}{\le d/2}$, $0\le i\le\tau$.
\paragraph{Lemma 6.34.} W.p. $>1-n^{-9\log n}$, for all $R\in\binom{[n]}{\le d/2}$,
\[\pm(Q_{0,0}^R-D_{\mathrm{Cl}}^\tau\mathbb E[Q_{0,0}^R]D_{\mathrm{Cl}}^\tau)(G)\preceq n^{-\epsilon}\operatorname{diag}(\widetilde{\mathrm{Cl}}(G))_{\mathcal S^R\times\mathcal S^R}.\tag{6.65}\]
\begin{proof}Fix $R$ and abbreviate $Q_{\mathrm{ps}}:=Q_{0,0}^R-D_{\mathrm{Cl}}^\tau\mathbb E[Q_{0,0}^R]D_{\mathrm{Cl}}^\tau=(Q_{\mathrm{ps},(i,j)})_{0\le i,j\le\tau}$. Block $(i,j)$ is supported on $S_{i,j}^2$, $S_{i,j}=\{A:|A|+\min\{i,j\}\ge d/2\}$. For $A\ne B$,
\begin{align*}Q_{\mathrm{ps},(i,j)}(A,B)&=Q_{0,0}^R((A,i),(B,j))\\&=\sum_{\substack{T_m:|V(T_m)\cup A\cup B|\le\tau\\A,B\in\mathrm{mSep}_{A,B}(T_m)}}(\omega/n)^{|V(T_m)\cup A\cup B|-(|A|+|B|)/2}q(A,B;T_m)\chi_{T_m}.\tag{6.66}\end{align*}
On the diagonal,
\[Q_{\mathrm{ps},(i,j)}(A,A)=\sum_{T_m:1\le|V(T_m)\setminus A|\le\tau-|A|}(\omega/n)^{|V(T_m)\cup A|-|A|}q(A,A;T_m)\chi_{T_m}.\tag{6.67}\]
Abbreviate $q(A,B;T_m):=Y_0(|R|,|V(T_m)\cup A\cup B|+(i+j))$. Two properties:
\[q(A,B;T_m)\text{ depends only on }|V(T_m)\cup A\cup B|\text{ for fixed }(A,B),\tag{6.68}\]
\[q(A,B;T_m)<\tau^{5\tau}\quad\text{by Lemma6.8(4)}.\tag{6.69}\]
Thus each entry factors through $\mathrm{Cl}_{A\cup B}$, hence $\mathrm{Cl}_A\mathrm{Cl}_B$. In particular,
\[Q_{\mathrm{ps}}=D'Q_{\mathrm{ps}}D'.\tag{6.70}\]
\paragraph{Claim 6.35.} W.p. $>1-n^{-9.5\log n}$, all blocks satisfy
\[\pm Q_{\mathrm{ps},(i,j)}\prec n^{-1.1\epsilon}\operatorname{diag}(2^{\binom{|A|}2})_{S^R_{\min\{i,j\}}\times S^R_{\min\{i,j\}}},\]
where $S_a^R=\{A\in\binom{[n]}{\le d/2}:|A|+a\ge d/2\}$.
The lemma follows from this claim and(6.70). For $b\in[0,d/2]$, let $I_b:=\{i:d/2-b\le i\le\tau\}$, $\mathcal P_b^R:=\{A\subseteq[n]:R\subseteq A,|A|=b\}$, and $Q_{\mathrm{ps};b}$ be the principal minor on $W_b:=(\mathcal P_b^R\times I_b)^2$, zero elsewhere. These $W_b$ partition the support where $|A|=|B|$. Remark6.16(2) gives
\[Q_{\mathrm{ps}}=\sum_{b=0}^{d/2}Q_{\mathrm{ps};b}.\tag{6.71}\]
Each is blockwise on $I_b^2$, blocks principal minors of $Q_{\mathrm{ps},(i,j)}$. By Claim6.35 w.p. $>1-n^{-9.5\log n}$ each signed block is $\prec n^{-1.5\epsilon}\operatorname{diag}(2^{\binom b2})$, so $\pm Q_{\mathrm{ps};b}\prec\tau^2n^{-1.5\epsilon}\operatorname{diag}(2^{\binom b2})_{W_b}\prec n^{-\epsilon}\operatorname{diag}(2^{\binom b2})_{W_b}$. By(6.71) and union-bound over $b$, $\pm Q_{\mathrm{ps}}\prec n^{-\epsilon}\operatorname{diag}(2^{\binom{|A|}2})_{\mathcal S^R\times\mathcal S^R}$ w.p. $1-n^{-9\log n}$. Insert this into(6.70), with $\widetilde{\mathrm{Cl}}_A=2^{\binom{|A|}2}\mathrm{Cl}_A$ and $\mathrm{Cl}_A^2=\mathrm{Cl}_A$, obtaining(6.65).\end{proof}
\paragraph{Proof of Claim 6.35.} Use Section4's norm bounds. Fix $(i,j)$, split $Q_{\mathrm{ps},(i,j)}$ into diagonal and off-diagonal parts.
\paragraph{Diagonal part.} For supported $(A,A)$, (6.67) gives
\[Q_{\mathrm{ps},(i,j)}^{\mathrm{diag}}(A,A)=\widetilde{\mathrm{Cl}}_A\underbrace{\sum_{\substack{T_m:1\le|V(T_m)\setminus A|\le\tau-|A|\\T_m\cap E[A]=\emptyset}}(\omega/n)^{|V(T_m)\setminus A|}q(A,A;T_m)\chi_{T_m}}_{=:g(A)}.\]
For fixed $A$, $q$ depends only on $|V(T_m)\setminus A|$, denote it $q(v)$. For $1\le v\le\tau-|A|$, let $U_1^v,\ldots,U_{h(v)}^v$ be all shapes $(A,A;T)$ with $T\cap E[A]=\emptyset$, $|V(T)\setminus A|=v$. Clearly,
\[h(v)\le2^{|A|v+v^2}.\tag{6.72}\]
W.p. $>1-n^{-9.6\log n}$,
\begin{align*}|g(A)|&=\left|\sum_{v=1}^{\tau-|A|}(\omega/n)^vq(v)\sum_{x=1}^{h(v)}\sum_{T_m:(A,A;T_m)\text{ has shape }U_x^v}\chi_{T_m}\right|\\&\le\sum_{v=1}^{\tau-|A|}(\omega/n)^vq(v)\sum_{x=1}^{h(v)}\|M_{U_x^v}\|\\&\le\sum_{v=1}^{\tau-|A|}(\omega/n)^v\tau^{5\tau}\sum_{x=1}^{h(v)}\|M_{U_x^v}\|\\&<\sum_{v=1}^{\tau}(\omega/n)^v\tau^{5\tau}2^{|A|v+v^2}n^{v/2}2^{O(|A|+v)}\\&<\sum_{v=1}^{\tau}n^{-3\epsilon v}n^{\epsilon v}<n^{-1.2\epsilon}.\end{align*}
The character sum equals $M_{U_x^v}(A,A)$ by Definition4.7; these matrices are diagonal. Use(6.69),(6.72),Theorem4.8 and the parameter regime in the respective inequalities.
\paragraph{Off-diagonal part.} By coefficient symmetry(6.68), it is a sum of graphical matrices. Let $U_1^{s,t},\ldots,U_{h(s,t)}^{s,t}$ be all distinct shapes $(A,B;T)$ with $|A|=|B|=s$, $A\ne B$, $A,B\in\mathrm{mSep}_{A,B}(T)$ and $|V(T)\cup A\cup B|=t$. By(6.66), it is block-diagonal by $s=d/2-i,\ldots,d/2$, with block
\[Q_{\mathrm{ps},(i,j)}^{\mathrm{off}}(s)=\sum_{s<t\le\tau}(\omega/n)^{t-s}\sum_{x=1}^{h(s,t)}q(U_x^{s,t})M_{U_x^{s,t}}.\]
By Theorem4.8 w.p. $>1-n^{-9.8\log n}$,
\[\|Q_{\mathrm{ps},(i,j)}^{\mathrm{off}}(s)\|\le\sum_{s<t\le\tau}(\omega/n)^{t-s}h(t,s)n^{(t-s)/2}2^{O(t)}(\log n)^{O(t-s)}.\tag{6.73}\]
Clearly $h(t,s)\le2^{\binom t2+O(t)}$. With the same probability,
\begin{align*}\text{RHS of(6.73)}&\le\sum_{\substack{d/2-\max\{i,j\}\le s\le d/2\\s<t\le\tau}}(\omega/n)^{t-s}2^{\binom t2+O(t)}n^{(t-s)/2}(\log n)^{O(t-s)}\\&<\sum_{\substack{d/2-\max\{i,j\}\le s\le d/2\\s<t\le\tau}}n^{-2\epsilon(t-s)}2^{O(t)}2^{\binom s2}(2^{t+s}\log n)^{O(t-s)}\\&<2^{\binom s2}n^{-1.9\epsilon},\end{align*}
by the parameter regime. Add these diagonal blocks to obtain $\pm Q_{\mathrm{ps},(i,j)}^{\mathrm{off}}\prec n^{-1.9\epsilon}\operatorname{diag}(2^{\binom{|A|}2})$ on the stated support. Finally union-bound: w.p. $>1-n^{-9.5\log n}$,
\[\pm Q_{\mathrm{ps},(i,j)}=\pm(Q_{\mathrm{ps},(i,j)}^{\mathrm{diag}}+Q_{\mathrm{ps},(i,j)}^{\mathrm{off}})\prec n^{-1.5\epsilon}\operatorname{diag}(2^{\binom{|A|}2})_{S^R_{\min\{i,j\}}\times S^R_{\min\{i,j\}}}.\]
This completes the claim.\qed
\paragraph{Corollary 6.36 (Positiveness of $Q_{0,0}^R$).} For every $R\in\binom{[n]}{\le d/2}$, w.p. $>1-n^{-8\log n}$,
\[Q_{0,0}^R(G)\succeq\tau^{-6.1\tau}\operatorname{diag}(\widetilde{\mathrm{Cl}}(G))_{\mathcal S^R\times\mathcal S^R}.\]
\begin{proof}By Lemma6.34 and Corollary6.32, where $\tau^{-6.1\tau}\gg n^{-\epsilon/10}$ in our parameter regime.\end{proof}

\section*{Author remaining block bounds and final PSD proof}
\paragraph{Lemma 6.37 (Bounds on $Q_{0,k}^R$).} W.p. $>1-n^{-9\log n}$, all $R\in\binom{[n]}{\le d/2}$ and $1\le k\le d/2$ satisfy
\[\pm Q_{0,k}^R(G)\preceq n^{-k/10}\operatorname{diag}(\widetilde{\mathrm{Cl}}(G))_{\mathcal S^R\times\mathcal S^R}.\]
\begin{proof}Union-bound over $(R,k)$, fixing them first, and abbreviate $Q=Q_{0,k}^R$, $q_k=q_{0,k}^R$. From Lemma6.19(3), $Q$ is supported on $\mathcal S^R\times\mathcal S^R$ and
\[Q((A,i),(B,j))=\sum_{T_m:|V(T_m)\cup A\cup B|\le\tau}(\omega/n)^{t-s}q_k(R_m,i,j)\chi_{T_m},\tag{6.74}\]
where $t=|A\cup B|$, $s=(|A|+|B|)/2$. For fixed $(i,j)$,
\[q_k(\cdot,i,j)\text{ is symmetric w.r.t. shapes},\tag{6.75}\]
\[|q_k(R_m,i,j)|\le\tau^{5\tau}(\omega/n^{1-\epsilon})^{s-p+k/3},\tag{6.76}\]
where $p$ is the max number of vertex-disjoint $A$--$B$ paths. Symmetry makes entries factor through $\mathrm{Cl}(A)\mathrm{Cl}(B)$, so
\[Q=D'QD'.\tag{6.77}\]
It suffices that, w.p. $>1-n^{-9.5\log n}$,
\[\pm Q\prec n^{-k/10}\operatorname{diag}(2^{\binom{|A|}2})_{\mathcal S^R\times\mathcal S^R}.\tag{6.78}\]
Insert this into(6.77) as in Lemma6.34. To prove(6.78), express blocks as graphical matrices. For fixed $(i,j)$ and $(s_1,s_2)\in\{0,\ldots,d/2\}^2$ with $s_1+i,s_2+j\ge d/2$, $t\ge\max\{s_1,s_2\}$, let $U_x^{t;s_1,s_2}$, $1\le x\le h(t;s_1,s_2)$, be all shapes with $|A|=s_1,|B|=s_2$, total size $t$. By(6.74) and symmetry,
\[Q_{(i,j)}=\sum_{\substack{(t;s_1,s_2)\colon s_1+i,s_2+j\ge d/2\\\tau\ge t\ge s_1,s_2}}\sum_{x=1}^{h(t;s_1,s_2)}q_k(U_x^{t;s_1,s_2},i,j)M_{U_x^{t;s_1,s_2}}.\]
View this blockwise by $(s_1,s_2)$:
\[Q_{(i,j)}=\sum_{s_1,s_2:\ s_1+i,s_2+j\ge d/2}Q_{(s_1,i),(s_2,j)},\tag{6.79}\]
where
\[Q_{(s_1,i),(s_2,j)}:=\sum_{s_1,s_2\le t\le\tau}\sum_{x=1}^{h(t;s_1,s_2)}q_k(U_x^{t;s_1,s_2},i,j)M_{U_x^{t;s_1,s_2}}.\tag{6.80}\]
This is a $\binom{[n]}{s_1}\times\binom{[n]}{s_2}$ matrix in block$(i,j)$. By Theorem4.8 and(6.76), w.p. $>1-n^{-10\log n}$,
\begin{multline*}\|Q_{(s_1,i),(s_2,j)}\|\le\sum_{\substack{t\le\tau\\t\ge s_1,s_2}}h(t;s_1,s_2)(\omega/n)^{t-s}(\omega/n^{1-\epsilon})^{s-p+k/3}\\\cdot n^{(t-p)/2}2^{O(t)}(\log n)^{O(t-s)},\tag{6.81}\end{multline*}
where $s=(s_1+s_2)/2$ and $p$ is the max path count in the shape. Since $h(t;s_1,s_2)\le2^{\binom t2+O(t)}=2^{\binom s2+O(t)+(t+s)(t-s)}$, and $k>0$, $2^{O(t)}<n^{\epsilon/10}$, $\tau^{5\tau}<n^{1/30}$, the RHS is
\[<2^{\binom s2}\tau^{5\tau}n^{-k/6}n^{-\epsilon(t-s)}<2^{\binom s2}n^{-k/8}.\tag{6.82}\]
Finally extend double-blocks by zero to $\mathcal S^R\times\mathcal S^R$ and sum:
\[Q=\sum_{(s_1,i),(s_2,j):s_1+i,s_2+j\ge d/2}Q_{(s_1,i),(s_2,j)}.\tag{6.83}\]
For each summand, (6.82) and Cauchy--Schwarz give
\[\pm Q_{(s_1,i),(s_2,j)}\prec\frac{n^{-k/8}}2\left(2^{\binom{s_1}2}\mathrm{Id}_{(s_1,i),(s_1,i)}+2^{\binom{s_2}2}\mathrm{Id}_{(s_2,j),(s_2,j)}\right).\]
Thus w.p. $>1-n^{-9.5\log n}$,
\[\pm Q\prec\tau^2n^{-k/8}\operatorname{diag}(2^{\binom{|A|}2})\prec n^{-k/10}\operatorname{diag}(2^{\binom{|A|}2})\]
on $\mathcal S^R\times\mathcal S^R$, proving(6.78).\end{proof}
\paragraph{Lemma 6.38 (Bounds for $c>0$).} W.p. $>1-n^{-9\log n}$, for every $R\in\binom{[n]}{\le d/2}$, $0<c\le|R|$, $0\le k\le d/2$,
\[\pm\omega^{-c}Q_{c,k}^R\preceq n^{-c/3}\operatorname{diag}(\widetilde{\mathrm{Cl}})_{\mathcal S^R\times\mathcal S^R}.\tag{6.84}\]
\begin{proof}Union-bound over $(R,c,k)$: it suffices to prove the inequality for a fixed triple with probability $>1-n^{-9.5\log n}$. If $k>0$, the proof is identical to Lemma6.37: by Lemma6.19 the same symmetry and coefficient-size conditions(6.75),(6.76) hold, with the same support. So consider $c>0,k=0$. Definition6.15 gives
\begin{align*}Q_{c,0}^R((A,i),(B,j))=\sum_{\substack{T_m:|V(T_m)\cup A\cup B|\le\tau\\A,B\in\mathrm{mSep}_{A,B}(T_m)}}&(\omega/n)^{|V(T_m)\cup A\cup B|-(|A|+|B|)/2}\\&\cdot Y_c(|R|,|V(T_m)\cup A\cup B|+(i+j))\chi_{T_m}.\tag{6.85}\end{align*}
Lemma6.8(4) bounds $Y_c$ by $\tau^{5\tau}$. Put $t=|V(T_m)\cup A\cup B|$, $s=(|A|+|B|)/2=|A|=|B|$. The coefficient is bounded by $(\omega/n)^{t-s}\tau^{5\tau}$. Support and symmetry let the proof proceed as before to(6.81), now with a single term
\[h(t;s_1,s_2)(\omega/n)^{t-s}\tau^{5\tau}n^{(t-p)/2}2^{O(t)}(\log n)^{O(t-s)}.\]
Every appearing ribbon has $A,B\in\mathrm{mSep}_{A,B}(T_m)$, so $p=s$ (the specialty of $k=0$). Replace(6.82)'s bound by $\tau^32^{\binom s2}n^{-3\epsilon(t-s)}\tau^{5\tau}2^{O(t)}<2^{\binom s2}\tau^{6\tau}$ and proceed to the last line as before, obtaining
\[\pm Q_{c,0}^R\prec\tau^{7\tau}\operatorname{diag}(2^{\binom{|A|}2})_{\mathcal S^R\times\mathcal S^R}.\]
Since $c\ge1$, $\omega=n^{1/2-4\epsilon}$ with $\epsilon<1/40$, and $\tau^{7\tau}<n^{1/15}$, parameters give
\[\pm\omega^{-c}Q_{c,0}^R\prec n^{-c/3}\operatorname{diag}(2^{\binom{|A|}2}).\]
Using again $Q_{c,0}^R=D'Q_{c,0}^RD'$ converts this to the normalized clique diagonal in(6.84).\end{proof}
Lemma6.21 now follows immediately from Corollary6.36 and Lemmas6.37/6.38.
\section*{Author Section 6.7: last step}
Recall $\mathcal P^R=\{I\in\binom{[n]}{d/2}:R\subseteq I\}$. Main Lemma6.9 requires, w.p. $1-n^{-5\log n}$ simultaneously,
\[M_0^R\succeq n^{-d}\operatorname{diag}(\widetilde{\mathrm{Cl}})_{\mathcal P^R\times\mathcal P^R},\tag{6.86}\]
\[\pm\omega^{-c}M_c^R\preceq n^{-c/6}M_0^R\quad(0<c\le|R|).\tag{6.87}\]
\paragraph{Lemma 6.39.} For every $R\in\binom{[n]}{\le d/2}$, on every $G$,
\[\widetilde L_RD^\tau\operatorname{diag}(\widetilde{\mathrm{Cl}})_{\mathcal S^R\times\mathcal S^R}D^\tau(\widetilde L_R)^\top\succeq(\omega/n)^{d/2}\operatorname{diag}(\widetilde{\mathrm{Cl}})_{\mathcal P^R\times\mathcal P^R}.\]
\begin{proof}Fix $R$. Regard supports on $[n']=[n]-R$ and reindex matrices by $\binom{[n']}{d'/2}$ or $\binom{[n']}{\le d'/2}$ with $d'/2=d/2-|R|$; $\tau$ is unchanged. Still write $\widetilde{\mathrm{Cl}}(X)$ for $\widetilde{\mathrm{Cl}}(X\sqcup R)$ for $X\subseteq[n']$. Since $D^\tau\operatorname{diag}(\widetilde{\mathrm{Cl}})D^\tau$ is nonnegative and diagonal,
\begin{multline*}\widetilde L_R[D^\tau\operatorname{diag}(\widetilde{\mathrm{Cl}})D^\tau](\widetilde L_R)^\top\\\succeq L_{R,0}[D^\tau\operatorname{diag}(\widetilde{\mathrm{Cl}})D^\tau](L_{R,0})^\top.\tag{6.88}\end{multline*}
Recall $\widetilde L_R=(L_{R,0},\ldots,L_{R,\tau})$. Further decompose $L_{R,0}=(L_{R,0}^0,\ldots,L_{R,0}^{d'/2})$ by the size of the column set. In particular,
\[L_{R,0}^{d/2-|R|}=(0,\ldots,0,\operatorname{diag}(\widetilde{\mathrm{Cl}})_{\binom{[n']}{d'/2}\times\binom{[n']}{d'/2}}),\]
since Definition6.11 allows only zero-reduced-size ribbons, which with its other conditions means $A=I$ and $T'\subseteq E(I)$. Thus the RHS of(6.88) is $\succ(\omega/n)^{d/2}\operatorname{diag}(\widetilde{\mathrm{Cl}})$ on that minor. Translating back gives the lemma.\end{proof}
\paragraph{Proof of Main Lemma 6.9.} Fix $R$. By Lemma6.19, for each $c\le|R|$,
\[M_c^R=\widetilde L_R[D^\tau(Q_{c,0}^R-Q_{c,1}^R+\cdots\pm Q_{c,d}^R)D^\tau](\widetilde L_R)^\top+E_c^R.\tag{6.89}\]
The following bounds each hold with probability $>1-n^{-8\log n}$, and union-bound makes the overall probability $>1-n^{-5\log n}$.
For(6.86),
\begin{align*}M_0^R&\succeq\tau^{-7\tau}\widetilde L_RD^\tau\operatorname{diag}(\widetilde{\mathrm{Cl}})D^\tau(\widetilde L_R)^\top+E_0^R\\&\qquad\text{Lemma6.21(1)}\\&\succeq\tau^{-7\tau}(\omega/n)^{d/2}\operatorname{diag}(\widetilde{\mathrm{Cl}})_{\mathcal P^R\times\mathcal P^R}+E_0^R\\&\qquad\text{Lemma6.39}\\&\succeq(\tau^{-7\tau}(\omega/n)^{d/2}-n^{-\epsilon\tau/2})\operatorname{diag}(\widetilde{\mathrm{Cl}})_{\mathcal P^R\times\mathcal P^R}\\&\qquad\text{Lemma6.19(4)}\\&\succeq n^{-d}\operatorname{diag}(\widetilde{\mathrm{Cl}})_{\mathcal P^R\times\mathcal P^R}\\&\qquad\text{parameter regime}.\end{align*}
For(6.87), fix $1\le c\le|R|$:
\begin{align*}M_c^R&\preceq\omega^cn^{-c/4}\widetilde L_RD^\tau\operatorname{diag}(\widetilde{\mathrm{Cl}})D^\tau(\widetilde L_R)^\top+E_c^R\\&\preceq\omega^cn^{-c/4}\tau^{7\tau}(M_0^R-E_0^R)+E_c^R\\&\preceq\omega^cn^{-c/5}M_0^R+(\omega^cn^{-c/5}+1)n^{-\epsilon\tau/2}\operatorname{diag}(\mathrm{Cl})_{\mathcal P^R\times\mathcal P^R}.\end{align*}
The respective steps use Lemma6.21(2), then Lemma6.21(1) and(6.89), then Lemma6.19(4). So
\begin{align*}\omega^{-c}M_c^R&\preceq n^{-c/5}M_0^R+2n^{-\epsilon\tau/2}\operatorname{diag}(\mathrm{Cl})_{\mathcal P^R\times\mathcal P^R}\\&\preceq(n^{-c/5}+2n^dn^{-\epsilon\tau/2})M_0^R\\&\qquad\text{by(6.86), }\widetilde{\mathrm{Cl}}\ge\mathrm{Cl}\\&\preceq n^{-c/6}M_0^R\\&\qquad\text{since }c\le|R|\le d/2\text{ and parameters}.\end{align*}
The same analysis holds for $-\omega^{-c}M_c^R$. This proves Main Lemma6.9, hence Theorem6.1 and the main exact-clique lower bound as shown above.\qed

\begin{thebibliography}{99}
\bibitem{ref1} Noga Alon, Michael Krivelevich, and Benny Sudakov. Finding a large hidden clique in a random graph. Random Structures \& Algorithms, 13(3-4):457–466, 1998.
\bibitem{ref2} Benny Applebaum, Boaz Barak, and Avi Wigderson. Public-key cryptography from different assumptions. In Proceedings of the forty-second ACM symposium on Theory of computing, pages 171–180, 2010.
\bibitem{ref3} Sanjeev Arora, Boaz Barak, Markus Brunnermeier, and Rong Ge. Computational complexity and information asymmetry in financial products. Communications of the ACM, 54(5):101–107, 2011.
\bibitem{ref4} Boaz Barak, Fernando GSL Brandao, Aram W Harrow, Jonathan Kelner, David Steurer, and Yuan Zhou. Hypercontractivity, sum-of-squares proofs, and their applications. In Proceedings of the forty-fourth annual ACM symposium on Theory of computing, pages 307–326, 2012.
\bibitem{ref5} Boaz Barak, Samuel Hopkins, Jonathan Kelner, Pravesh K Kothari, Ankur Moitra, and Aaron Potechin. A nearly tight sum-of-squares lower bound for the planted clique problem. SIAM Journal on Computing, 48(2):687–735, 2019.
\bibitem{ref6} Boaz Barak and David Steurer. Sum-of-squares proofs and the quest toward optimal algorithms. In Proceedings of International Congress of Mathematicians (ICM), 2014.
\bibitem{ref7} Paul Beame, Russell Impagliazzo, Jan Krajíček, Toniann Pitassi, and Pavel Pudlák. Lower bounds on hilbert’s nullstellensatz and propositional proofs. Proceedings of the London Mathematical Society, 3(1):1–26, 1996.
\bibitem{ref8} Quentin Berthet and Philippe Rigollet. Complexity theoretic lower bounds for sparse principal component detection. In Conference on Learning Theory, pages 1046–1066. PMLR, 2013.
\bibitem{ref9} P Delsarte. An algebraic approach to association schemes of coding theory, phillips j, 1973.
\bibitem{ref10} Yash Deshpande and Andrea Montanari. Improved sum-of-squares lower bounds for hidden clique and hidden submatrix problems. In Conference on Learning Theory, pages 523–562. PMLR, 2015.
\bibitem{ref11} Fernando Escalante. Schnittverbände in graphen. In Abhandlungen aus dem Mathematischen Seminar der Universität Hamburg, volume 38, pages 199–220. Springer, 1972.
\bibitem{ref12} Uriel Feige and Robert Krauthgamer. Finding and certifying a large hidden clique in a semirandom graph. Random Structures \& Algorithms, 16(2):195–208, 2000.
\bibitem{ref13} Uriel Feige and Robert Krauthgamer. The probable value of the lovász–schrijver relaxations for maximum independent set. SIAM Journal on Computing, 32(2):345–370, 2003.
\bibitem{ref14} Dima Grigoriev and Nicolai Vorobjov. Complexity of null-and positivstellensatz proofs. Annals of Pure and Applied Logic, 113(1-3):153–160, 2001.
\bibitem{ref15} Samuel B Hopkins, Pravesh Kothari, Aaron Henry Potechin, Prasad Raghavendra, and Tselil Schramm. On the integrality gap of degree-4 sum of squares for planted clique. ACM Transactions on Algorithms (TALG), 14(3):1–31, 2018.
\bibitem{ref16} Samuel B Hopkins, Pravesh K Kothari, and Aaron Potechin. Sos and planted clique: Tight analysis of mpw moments at all degrees and an optimal lower bound at degree four. arXiv preprint, 2015. arXiv:1507.05230.
\bibitem{ref17} Samuel B Hopkins, Pravesh K Kothari, Aaron Potechin, Prasad Raghavendra, Tselil Schramm, and David Steurer. The power of sum-of-squares for detecting hidden structures. In 2017 IEEE 58th Annual Symposium on Foundations of Computer Science (FOCS), pages 720–731. IEEE, 2017.
\bibitem{ref18} Mark Jerrum. Large cliques elude the metropolis process. Random Structures \& Algorithms, 3(4):347–359, 1992.
\bibitem{ref19} R Karp. Probabilistic analysis of some combinatorial search problems. traub, jf (ed.): Algo- rithms and complexity: New directions and recent results, 1976.
\bibitem{ref20} Pravesh Kothari, Ryan O’Donnell, and Tselil Schramm. Sos lower bounds with hard constraints: think global, act local. arXiv preprint, 2018. arXiv:1809.01207.
\bibitem{ref21} Pravesh K Kothari and Ruta Mehta. Sum-of-squares meets nash: Optimal lower bounds for finding any equilibrium. arXiv preprint, 2018. arXiv:1806.09426.
\bibitem{ref22} Luděk Kučera. Expected complexity of graph partitioning problems. Discrete Applied Mathematics, 57(2-3):193–212, 1995.
\bibitem{ref23} Jean B Lasserre. Global optimization with polynomials and the problem of moments. SIAM Journal on optimization, 11(3):796–817, 2001.
\bibitem{ref24} Dhruv Medarametla and Aaron Potechin. Bounds on the norms of uniform low degree graph matrices. In Approximation, Randomization, and Combinatorial Optimization. Algorithms and Techniques (APPROX/RANDOM 2016). Schloss Dagstuhl-Leibniz-Zentrum fuer Informatik, 2016.
\bibitem{ref25} Raghu Meka, Aaron Potechin, and Avi Wigderson. Sum-of-squares lower bounds for planted clique. In Proceedings of the forty-seventh annual ACM symposium on Theory of computing, pages 87–96, 2015.
\bibitem{ref26} Ryan O’Donnell. Sos is not obviously automatizable, even approximately. In 8th Innovations in Theoretical Computer Science Conference (ITCS 2017). Schloss Dagstuhl-Leibniz-Zentrum fuer Informatik, 2017.
\bibitem{ref27} Pablo A Parrilo. Structured semidefinite programs and semialgebraic geometry methods in robustness and optimization. PhD thesis, California Institute of Technology, 2000.
\bibitem{ref28} Pavel A Pevzner, Sing-Hoi Sze, et al. Combinatorial approaches to finding subtle signals in dna sequences. In ISMB, volume 8, pages 269–278, 2000.
\bibitem{ref29} Prasad Raghavendra and Benjamin Weitz. On the bit complexity of sum-of-squares proofs. arXiv preprint, 2017. arXiv:1702.05139.
\bibitem{ref30} Naum Z Shor. Class of global minimum bounds of polynomial functions. Cybernetics, 23(6):731–734, 1987.
\bibitem{ref31} Evgenij E Tyrtyshnikov. How bad are hankel matrices? Numerische Mathematik, 67(2):261– 269, 1994.
\end{thebibliography}
\end{document}
